\documentclass[8pt]{beamer}
\usetheme{Madrid}
\usecolortheme{default}
\usepackage{amsmath, amssymb, graphicx, array, multirow, booktabs, xcolor, tikz}
\usetikzlibrary{shapes,arrows,positioning,fit,backgrounds}

% Compact font settings
\setbeamerfont{title}{size=\Large}
\setbeamerfont{subtitle}{size=\small}
\setbeamerfont{author}{size=\scriptsize}
\setbeamerfont{date}{size=\scriptsize}
\setbeamerfont{frametitle}{size=\large}
\setbeamerfont{framesubtitle}{size=\normalsize}
\setbeamerfont{enumerate item}{size=\footnotesize}
\setbeamerfont{itemize item}{size=\footnotesize}
\setbeamerfont{block title}{size=\footnotesize}

% Compact spacing
\setlength{\itemsep}{0.08ex}
\setlength{\parskip}{0.08ex}
\setlength{\leftmargini}{0.3cm}
\setbeamersize{text margin left=3mm, text margin right=3mm}
\addtobeamertemplate{frame begin}{\vspace{-0.4cm}}{}

% Colors
\definecolor{myblue}{RGB}{0,0,255}
\definecolor{myred}{RGB}{255,0,0}
\definecolor{mygreen}{RGB}{0,128,0}
\definecolor{myorange}{RGB}{255,165,0}
\definecolor{mypurple}{RGB}{128,0,128}
\definecolor{mygray}{RGB}{128,128,128}

\title{Supervised Learning - Model Estimation}
\subtitle{100 Questions: OLS, MLE, Gradient Descent, Regularization, Cross-Validation}
\author{GARP RAI Exam Preparation}
\date{}

\begin{document}
	
	\begin{frame}
		\titlepage
	\end{frame}
	
	% ============================================================
	% SECTION 1: ESTIMATION PARADIGMS - QUESTIONS 1-15
	% ============================================================
	
	\begin{frame}{Q1: OLS Definition}
		\fontsize{6.5}{7.5}\selectfont
		\textbf{What is Ordinary Least Squares (OLS) estimation in supervised learning?}
		
		\begin{enumerate}
			\item A method that minimizes the sum of absolute errors between predicted and actual values.
			\item A method that minimizes the sum of squared errors between predicted and actual values.
			\item A method that maximizes the likelihood of the observed data.
			\item A method that minimizes the number of parameters in the model.
		\end{enumerate}
		
		\vspace{0.05cm}
		\textcolor{myblue}{\textbf{Answer: B}}
		\vspace{0.05cm}
		\textcolor{myblue}{\textbf{Explanation:}} OLS minimizes $\sum_{i=1}^N (y_i - \hat{y}_i)^2$ - the sum of squared residuals.
		
		\vspace{0.05cm}
		\textcolor{myred}{\textbf{Why Others Are Wrong:}}
		\begin{itemize}
			\item \textcolor{myred}{A:} This describes Least Absolute Deviations (LAD) or L1 loss - a different estimator that minimizes absolute errors.
			\item \textcolor{myred}{C:} This describes Maximum Likelihood Estimation (MLE) - a different estimation paradigm that maximizes likelihood rather than minimizing errors.
			\item \textcolor{myred}{D:} This describes regularization (specifically L1/Lasso) which penalizes model complexity - not OLS.
		\end{itemize}
		
		\textcolor{myred}{\textbf{Exam Trap:}} OLS minimizes \textcolor{myblue}{squared errors}, NOT absolute errors! This is the most common confusion on exams.
	\end{frame}
	
	\begin{frame}{Q2: OLS Objective Function}
		\fontsize{6.5}{7.5}\selectfont
		\textbf{Which of the following is the correct objective function for Ordinary Least Squares?}
		
		\begin{enumerate}
			\item $L = \sum_{i=1}^N |y_i - \hat{y}_i|$
			\item $L = \sum_{i=1}^N (y_i - \hat{y}_i)^2$
			\item $L = \sum_{i=1}^N (y_i - \hat{y}_i)^4$
			\item $L = \max_i |y_i - \hat{y}_i|$
		\end{enumerate}
		
		\vspace{0.05cm}
		\textcolor{myblue}{\textbf{Answer: B}}
		\vspace{0.05cm}
		\textcolor{myblue}{\textbf{Explanation:}} OLS minimizes the Residual Sum of Squares (RSS): $\sum (y_i - \hat{y}_i)^2$.
		
		\vspace{0.05cm}
		\textcolor{myred}{\textbf{Why Others Are Wrong:}}
		\begin{itemize}
			\item \textcolor{myred}{A:} This is L1 loss (Least Absolute Deviations) - minimizes absolute errors, not squared errors.
			\item \textcolor{myred}{C:} Higher-order loss (L4) - rarely used in practice; doesn't correspond to OLS.
			\item \textcolor{myred}{D:} Chebyshev loss (minimax) - minimizes maximum absolute error, not sum of squared errors.
		\end{itemize}
		
		\textcolor{myred}{\textbf{Exam Trap:}} OLS = \textcolor{myblue}{RSS} = $\sum (y_i - \hat{y}_i)^2$! Don't confuse with MAE or other loss functions. The square term is the key identifier.
	\end{frame}
	
	\begin{frame}{Q3: OLS vs NLS}
		\fontsize{6.5}{7.5}\selectfont
		\textbf{What is the key difference between Ordinary Least Squares (OLS) and Nonlinear Least Squares (NLS)?}
		
		\begin{enumerate}
			\item OLS minimizes absolute errors while NLS minimizes squared errors.
			\item OLS is used when the model is linear in parameters; NLS is used when the model is nonlinear in parameters.
			\item OLS requires gradient descent while NLS has a closed-form solution.
			\item OLS is for classification while NLS is for regression.
		\end{enumerate}
		
		\vspace{0.05cm}
		\textcolor{myblue}{\textbf{Answer: B}}
		\vspace{0.05cm}
		\textcolor{myblue}{\textbf{Explanation:}} OLS requires linearity in parameters (e.g., $y = \beta_0 + \beta_1 x$). NLS handles nonlinearity in parameters (e.g., $y = \beta_0 + \beta_1 e^{\beta_2 x}$).
		
		\vspace{0.05cm}
		\textcolor{myred}{\textbf{Why Others Are Wrong:}}
		\begin{itemize}
			\item \textcolor{myred}{A:} Both OLS and NLS minimize squared errors - this is a fundamental misunderstanding of both methods.
			\item \textcolor{myred}{C:} Actually reversed! OLS has closed-form solution $\hat{\beta} = (X^TX)^{-1}X^Ty$; NLS requires iterative methods.
			\item \textcolor{myred}{D:} Both are regression methods (continuous dependent variables), not classification.
		\end{itemize}
		
		\textcolor{myred}{\textbf{Exam Trap:}} The difference is \textcolor{myblue}{linearity in parameters}, NOT in variables! $y = \beta_0 + \beta_1 x^2$ is OLS because it's linear in parameters, even though it's nonlinear in x.
	\end{frame}
	
	\begin{frame}{Q4: Nonlinear Least Squares - Example}
		\fontsize{6.5}{7.5}\selectfont
		\textbf{Which of the following models would require Nonlinear Least Squares (NLS) estimation?}
		
		\begin{enumerate}
			\item $y = \beta_0 + \beta_1 x_1 + \beta_2 x_2 + \varepsilon$
			\item $y = \beta_0 + \beta_1 x + \beta_2 x^2 + \varepsilon$
			\item $y = \beta_0 + \beta_1 e^{\beta_2 x} + \varepsilon$
			\item $y = \beta_0 + \beta_1 \ln(x) + \varepsilon$
		\end{enumerate}
		
		\vspace{0.05cm}
		\textcolor{myblue}{\textbf{Answer: C}}
		\vspace{0.05cm}
		\textcolor{myblue}{\textbf{Explanation:}} Option C has $\beta_1 e^{\beta_2 x}$ which is nonlinear in parameter $\beta_2$.
		
		\vspace{0.05cm}
		\textcolor{myred}{\textbf{Why Others Are Wrong:}}
		\begin{itemize}
			\item \textcolor{myred}{A:} This is a standard linear regression - linear in parameters $\beta_0, \beta_1, \beta_2$.
			\item \textcolor{myred}{B:} Polynomial regression - linear in parameters despite being nonlinear in x.
			\item \textcolor{myred}{D:} Log transformation - linear in parameters (treat $\ln(x)$ as transformed feature).
		\end{itemize}
		
		\textcolor{myred}{\textbf{Exam Trap:}} NLS is about \textcolor{myblue}{nonlinearity in parameters}, not variables! This is the single most common mistake on this topic.
	\end{frame}
	
	\begin{frame}{Q5: Maximum Likelihood Estimation - Definition}
		\fontsize{6.5}{7.5}\selectfont
		\textbf{What is Maximum Likelihood Estimation (MLE)?}
		
		\begin{enumerate}
			\item A method that minimizes the sum of squared errors.
			\item A method that finds the parameter values that maximize the probability (likelihood) of observing the given data.
			\item A method that minimizes the number of parameters in the model.
			\item A method that finds the parameter values that minimize the variance of the errors.
		\end{enumerate}
		
		\vspace{0.05cm}
		\textcolor{myblue}{\textbf{Answer: B}}
		\vspace{0.05cm}
		\textcolor{myblue}{\textbf{Explanation:}} MLE finds $\theta$ that maximizes $L(\theta) = P(\text{data}|\theta)$. It asks: "What parameters make the observed data most probable?"
		
		\vspace{0.05cm}
		\textcolor{myred}{\textbf{Why Others Are Wrong:}}
		\begin{itemize}
			\item \textcolor{myred}{A:} This describes OLS (sum of squared errors) - a completely different estimation paradigm.
			\item \textcolor{myred}{C:} This describes regularization (penalizing complexity) - not MLE.
			\item \textcolor{myred}{D:} Minimizing error variance is related to efficiency but is not the definition of MLE.
		\end{itemize}
		
		\textcolor{myred}{\textbf{Exam Trap:}} MLE maximizes \textcolor{myblue}{likelihood}, NOT minimizes error! This is the fundamental difference from OLS. Also, likelihood = $P(\text{data}|\theta)$, NOT $P(\theta|\text{data})$ (that's Bayesian posterior).
	\end{frame}
	
	\begin{frame}{Q6: MLE vs OLS - Normal Errors}
		\fontsize{6.5}{7.5}\selectfont
		\textbf{Under what condition do OLS and MLE produce the same parameter estimates for a linear regression model?}
		
		\begin{enumerate}
			\item When the errors are independent.
			\item When the errors are normally distributed with constant variance.
			\item When the errors are uniformly distributed.
			\item When the errors are autocorrelated.
		\end{enumerate}
		
		\vspace{0.05cm}
		\textcolor{myblue}{\textbf{Answer: B}}
		\vspace{0.05cm}
		\textcolor{myblue}{\textbf{Explanation:}} When $\varepsilon \sim N(0, \sigma^2)$, log-likelihood is proportional to negative RSS. Maximizing likelihood = minimizing RSS.
		
		\vspace{0.05cm}
		\textcolor{myred}{\textbf{Why Others Are Wrong:}}
		\begin{itemize}
			\item \textcolor{myred}{A:} Independence alone is insufficient - we need the specific normal distribution form.
			\item \textcolor{myred}{C:} Uniform errors would give a different likelihood function (MLE would minimize maximum absolute error).
			\item \textcolor{myred}{D:} Autocorrelation violates OLS assumptions and would not lead to OLS-MLE equivalence.
		\end{itemize}
		
		\textcolor{myred}{\textbf{Exam Trap:}} OLS and MLE are \textcolor{myblue}{equivalent} under normal errors! This is a critical exam concept - not just independence, but specifically normal distribution.
	\end{frame}
	
	\begin{frame}{Q7: MLE - Logistic Regression}
		\fontsize{6.5}{7.5}\selectfont
		\textbf{What estimation method is typically used for logistic regression?}
		
		\begin{enumerate}
			\item Ordinary Least Squares (OLS)
			\item Nonlinear Least Squares (NLS)
			\item Maximum Likelihood Estimation (MLE)
			\item Method of Moments
		\end{enumerate}
		
		\vspace{0.05cm}
		\textcolor{myblue}{\textbf{Answer: C}}
		\vspace{0.05cm}
		\textcolor{myblue}{\textbf{Explanation:}} Logistic regression uses MLE with Bernoulli likelihood. The log-likelihood is $\sum [y_i \log(p_i) + (1-y_i)\log(1-p_i)]$.
		
		\vspace{0.05cm}
		\textcolor{myred}{\textbf{Why Others Are Wrong:}}
		\begin{itemize}
			\item \textcolor{myred}{A:} OLS inappropriate because target is binary (0/1) and errors are not normal. OLS would produce heteroscedastic errors and predictions outside [0,1].
			\item \textcolor{myred}{B:} NLS is not the standard approach; logistic regression's nonlinearity is in the link function, handled naturally by MLE.
			\item \textcolor{myred}{D:} Method of Moments is less common and less efficient for logistic regression.
		\end{itemize}
		
		\textcolor{myred}{\textbf{Exam Trap:}} Logistic regression uses \textcolor{myblue}{MLE}, NOT OLS! This is a critical distinction for classification problems.
	\end{frame}
	
	\begin{frame}{Q8: Likelihood Function}
		\fontsize{6.5}{7.5}\selectfont
		\textbf{What does the likelihood function $L(\theta | \text{data})$ represent?}
		
		\begin{enumerate}
			\item The probability of the parameters given the data.
			\item The probability of the data given the parameters.
			\item The sum of squared errors for a given parameter set.
			\item The variance of the parameter estimates.
		\end{enumerate}
		
		\vspace{0.05cm}
		\textcolor{myblue}{\textbf{Answer: B}}
		\vspace{0.05cm}
		\textcolor{myblue}{\textbf{Explanation:}} $L(\theta | \text{data}) = P(\text{data} | \theta)$. It's the probability of observing the data given the parameters.
		
		\vspace{0.05cm}
		\textcolor{myred}{\textbf{Why Others Are Wrong:}}
		\begin{itemize}
			\item \textcolor{myred}{A:} This describes the posterior distribution in Bayesian inference ($P(\theta|\text{data})$).
			\item \textcolor{myred}{C:} This is the OLS loss function (RSS) - a different concept.
			\item \textcolor{myred}{D:} This describes the variance of estimators - a different statistical concept.
		\end{itemize}
		
		\textcolor{myred}{\textbf{Exam Trap:}} Likelihood = $P(\text{data} | \theta)$, NOT $P(\theta | \text{data})$! This confusion between likelihood and posterior is extremely common on exams.
	\end{frame}
	
	\begin{frame}{Q9: Log-Likelihood}
		\fontsize{6.5}{7.5}\selectfont
		\textbf{Why is the log-likelihood often used instead of the likelihood function?}
		
		\begin{enumerate}
			\item Because the log-likelihood is always positive.
			\item Because the log-likelihood is easier to maximize (it turns products into sums and is monotonic).
			\item Because the log-likelihood has a closed-form solution.
			\item Because the log-likelihood minimizes the number of parameters.
		\end{enumerate}
		
		\vspace{0.05cm}
		\textcolor{myblue}{\textbf{Answer: B}}
		\vspace{0.05cm}
		\textcolor{myblue}{\textbf{Explanation:}} Log transforms products to sums (easier to differentiate). Since log is monotonic, maximizing log-likelihood $\equiv$ maximizing likelihood.
		
		\vspace{0.05cm}
		\textcolor{myred}{\textbf{Why Others Are Wrong:}}
		\begin{itemize}
			\item \textcolor{myred}{A:} False - log-likelihood is often negative (since probabilities are $<1$, logs are negative).
			\item \textcolor{myred}{C:} False - log-likelihood usually requires numerical optimization (except simple cases).
			\item \textcolor{myred}{D:} False - log-likelihood doesn't affect parameter count; that's regularization.
		\end{itemize}
		
		\textcolor{myred}{\textbf{Exam Trap:}} Log-likelihood is \textcolor{myblue}{monotonic} - maximizing it is equivalent to maximizing likelihood! This is why we use it despite often being negative.
	\end{frame}
	
	\begin{frame}{Q10: Closed-Form vs Numerical Solutions}
		\fontsize{6.5}{7.5}\selectfont
		\textbf{Which estimation method typically has a closed-form (analytical) solution?}
		
		\begin{enumerate}
			\item Maximum Likelihood Estimation for logistic regression
			\item Nonlinear Least Squares
			\item Ordinary Least Squares for linear regression
			\item Maximum Likelihood Estimation for neural networks
		\end{enumerate}
		
		\vspace{0.05cm}
		\textcolor{myblue}{\textbf{Answer: C}}
		\vspace{0.05cm}
		\textcolor{myblue}{\textbf{Explanation:}} OLS for linear regression has closed-form: $\hat{\beta} = (X^T X)^{-1} X^T y$.
		
		\vspace{0.05cm}
		\textcolor{myred}{\textbf{Why Others Are Wrong:}}
		\begin{itemize}
			\item \textcolor{myred}{A:} Logistic regression MLE requires numerical optimization (Newton-Raphson, IRLS).
			\item \textcolor{myred}{B:} NLS requires iterative numerical optimization (no closed-form).
			\item \textcolor{myred}{D:} Neural networks use backpropagation (numerical gradient-based optimization).
		\end{itemize}
		
		\textcolor{myred}{\textbf{Exam Trap:}} Only \textcolor{myblue}{OLS for linear regression} has a closed-form solution! All other methods require numerical optimization. This is a key distinguishing feature.
	\end{frame}
	
	\begin{frame}{Q11: OLS Assumptions}
		\fontsize{6.5}{7.5}\selectfont
		\textbf{Which of the following is NOT a standard assumption of Ordinary Least Squares?}
		
		\begin{enumerate}
			\item Linearity in parameters
			\item Homoscedasticity (constant variance of errors)
			\item Normality of errors
			\item The errors are autocorrelated
		\end{enumerate}
		
		\vspace{0.05cm}
		\textcolor{myblue}{\textbf{Answer: D}}
		\vspace{0.05cm}
		\textcolor{myblue}{\textbf{Explanation:}} OLS assumes: linearity in parameters, homoscedasticity, independence (no autocorrelation), and (for inference) normality.
		
		\vspace{0.05cm}
		\textcolor{myred}{\textbf{Why Others Are Wrong:}}
		\begin{itemize}
			\item \textcolor{myred}{A:} This IS an OLS assumption - model must be linear in parameters.
			\item \textcolor{myred}{B:} This IS an OLS assumption - constant error variance (homoscedasticity).
			\item \textcolor{myred}{C:} This IS an OLS assumption (for hypothesis testing and confidence intervals).
		\end{itemize}
		
		\textcolor{myred}{\textbf{Exam Trap:}} OLS assumes \textcolor{myblue}{no autocorrelation} in errors! Autocorrelation violates the independence assumption. Option D is the violation, not the assumption.
	\end{frame}
	
	\begin{frame}{Q12: MLE - Properties}
		\fontsize{6.5}{7.5}\selectfont
		\textbf{Which of the following is a desirable property of Maximum Likelihood Estimators?}
		
		\begin{enumerate}
			\item They are always unbiased.
			\item They are consistent (converge to true values as sample size increases).
			\item They always have the smallest variance among all estimators.
			\item They never require numerical optimization.
		\end{enumerate}
		
		\vspace{0.05cm}
		\textcolor{myblue}{\textbf{Answer: B}}
		\vspace{0.05cm}
		\textcolor{myblue}{\textbf{Explanation:}} MLE is consistent: $\hat{\theta}_{MLE} \xrightarrow{p} \theta_0$ as $n \to \infty$ (under regularity conditions).
		
		\vspace{0.05cm}
		\textcolor{myred}{\textbf{Why Others Are Wrong:}}
		\begin{itemize}
			\item \textcolor{myred}{A:} False - MLE can be biased (e.g., variance estimator $\hat{\sigma}^2_{MLE}$ is biased).
			\item \textcolor{myred}{C:} False - MLE is asymptotically efficient but not always minimum variance in finite samples.
			\item \textcolor{myred}{D:} False - MLE often requires numerical optimization (logistic regression, etc.).
		\end{itemize}
		
		\textcolor{myred}{\textbf{Exam Trap:}} MLE is \textcolor{myblue}{consistent} but NOT necessarily unbiased! This is a common misconception - consistency and unbiasedness are different properties.
	\end{frame}
	
	\begin{frame}{Q13: NLS - Optimization}
		\fontsize{6.5}{7.5}\selectfont
		\textbf{Why does Nonlinear Least Squares typically require iterative numerical optimization?}
		
		\begin{enumerate}
			\item Because the objective function has multiple local minima.
			\item Because there is no closed-form solution for nonlinear-in-parameters models.
			\item Because the errors are not normally distributed.
			\item Because the model has too many parameters.
		\end{enumerate}
		
		\vspace{0.05cm}
		\textcolor{myblue}{\textbf{Answer: B}}
		\vspace{0.05cm}
		\textcolor{myblue}{\textbf{Explanation:}} NLS requires numerical optimization because the normal equations are nonlinear in parameters and have no closed-form solution.
		
		\vspace{0.05cm}
		\textcolor{myred}{\textbf{Why Others Are Wrong:}}
		\begin{itemize}
			\item \textcolor{myred}{A:} Multiple local minima are a potential issue but not the primary reason for numerical methods.
			\item \textcolor{myred}{C:} Error distribution doesn't determine the need for numerical optimization.
			\item \textcolor{myred}{D:} Parameter count affects complexity but isn't the fundamental reason.
		\end{itemize}
		
		\textcolor{myred}{\textbf{Exam Trap:}} No closed-form solution exists for \textcolor{myblue}{nonlinear-in-parameters} models! This is the fundamental reason for iterative methods, not local minima or parameter count.
	\end{frame}
	
	\begin{frame}{Q14: OLS - Gauss-Markov Theorem}
		\fontsize{6.5}{7.5}\selectfont
		\textbf{What does the Gauss-Markov theorem state about OLS estimators?}
		
		\begin{enumerate}
			\item OLS estimators are always unbiased.
			\item OLS estimators are the Best Linear Unbiased Estimators (BLUE) under certain conditions.
			\item OLS estimators always have the smallest variance among all estimators.
			\item OLS estimators are consistent for any distribution of errors.
		\end{enumerate}
		
		\vspace{0.05cm}
		\textcolor{myblue}{\textbf{Answer: B}}
		\vspace{0.05cm}
		\textcolor{myblue}{\textbf{Explanation:}} Under linearity, no perfect multicollinearity, zero conditional mean, homoscedasticity, OLS is BLUE.
		
		\vspace{0.05cm}
		\textcolor{myred}{\textbf{Why Others Are Wrong:}}
		\begin{itemize}
			\item \textcolor{myred}{A:} Close but incomplete - misses the "Best" (minimum variance) part.
			\item \textcolor{myred}{C:} False - only among linear unbiased estimators, not all estimators.
			\item \textcolor{myred}{D:} False - consistency requires additional conditions beyond Gauss-Markov.
		\end{itemize}
		
		\textcolor{myred}{\textbf{Exam Trap:}} OLS is \textcolor{myblue}{BLUE} under Gauss-Markov assumptions! Not "always best" - only among linear unbiased estimators. Non-linear estimators could have lower variance.
	\end{frame}
	
	\begin{frame}{Q15: Estimation Methods - Summary}
		\fontsize{6.5}{7.5}\selectfont
		\textbf{Which estimation method would you use for a model where $y = \beta_0 + \beta_1 x_1^{\beta_2} + \varepsilon$?}
		
		\begin{enumerate}
			\item Ordinary Least Squares (OLS)
			\item Nonlinear Least Squares (NLS)
			\item Maximum Likelihood Estimation (MLE)
			\item Both NLS and MLE would be appropriate
		\end{enumerate}
		
		\vspace{0.05cm}
		\textcolor{myblue}{\textbf{Answer: D}}
		\vspace{0.05cm}
		\textcolor{myblue}{\textbf{Explanation:}} Model is nonlinear in parameter $\beta_2$. Both NLS (minimizes squared errors) and MLE (maximizes likelihood) can be used.
		
		\vspace{0.05cm}
		\textcolor{myred}{\textbf{Why Others Are Wrong:}}
		\begin{itemize}
			\item \textcolor{myred}{A:} OLS requires linearity in parameters - this model violates that.
			\item \textcolor{myred}{B:} NLS alone is correct but incomplete - MLE also works.
			\item \textcolor{myred}{C:} MLE alone is correct but incomplete - NLS also works.
		\end{itemize}
		
		\textcolor{myred}{\textbf{Exam Trap:}} For nonlinear models, both \textcolor{myblue}{NLS and MLE} can be used! Choice depends on error distribution assumptions. If errors are normal, they're equivalent.
	\end{frame}
	
	% ============================================================
	% SECTION 2: GRADIENT DESCENT - QUESTIONS 16-35
	% ============================================================
	
	\begin{frame}{Q16: Gradient Descent - Definition}
		\fontsize{6.5}{7.5}\selectfont
		\textbf{What is gradient descent in the context of model estimation?}
		
		\begin{enumerate}
			\item A method that finds the global minimum of any function.
			\item An iterative optimization algorithm that updates parameters in the direction of the negative gradient of the loss function.
			\item A method that solves for parameters using a closed-form equation.
			\item A method that randomly searches for optimal parameters.
		\end{enumerate}
		
		\vspace{0.05cm}
		\textcolor{myblue}{\textbf{Answer: B}}
		\vspace{0.05cm}
		\textcolor{myblue}{\textbf{Explanation:}} Update rule: $\theta_{new} = \theta_{old} - \alpha \nabla L(\theta)$. It's iterative and uses the negative gradient direction.
		
		\vspace{0.05cm}
		\textcolor{myred}{\textbf{Why Others Are Wrong:}}
		\begin{itemize}
			\item \textcolor{myred}{A:} False - gradient descent can get stuck in local minima; no guarantee of global minimum.
			\item \textcolor{myred}{C:} This describes closed-form solutions (like OLS), not gradient descent.
			\item \textcolor{myred}{D:} This describes random search or Monte Carlo methods, not gradient descent.
		\end{itemize}
		
		\textcolor{myred}{\textbf{Exam Trap:}} Gradient descent is \textcolor{myblue}{iterative} and uses the \textcolor{myblue}{negative gradient}! Not a direct solution method, and doesn't guarantee global minimum.
	\end{frame}
	
	\begin{frame}{Q17: Gradient Descent Update Rule}
		\fontsize{6.5}{7.5}\selectfont
		\textbf{What is the general update rule for gradient descent?}
		
		\begin{enumerate}
			\item $\theta_{new} = \theta_{old} + \alpha \nabla L(\theta)$
			\item $\theta_{new} = \theta_{old} - \alpha \nabla L(\theta)$
			\item $\theta_{new} = \theta_{old} \times \alpha \nabla L(\theta)$
			\item $\theta_{new} = \theta_{old} / \alpha \nabla L(\theta)$
		\end{enumerate}
		
		\vspace{0.05cm}
		\textcolor{myblue}{\textbf{Answer: B}}
		\vspace{0.05cm}
		\textcolor{myblue}{\textbf{Explanation:}} $\theta_{new} = \theta_{old} - \alpha \nabla L(\theta)$. We subtract the gradient to move in the direction of decreasing loss.
		
		\vspace{0.05cm}
		\textcolor{myred}{\textbf{Why Others Are Wrong:}}
		\begin{itemize}
			\item \textcolor{myred}{A:} This is gradient ascent (moving to increase loss) - would maximize rather than minimize.
			\item \textcolor{myred}{C:} Multiplication is not the correct operation; gradient descent uses subtraction.
			\item \textcolor{myred}{D:} Division is not the correct operation; this is not a valid update rule.
		\end{itemize}
		
		\textcolor{myred}{\textbf{Exam Trap:}} Gradient descent uses \textcolor{myblue}{subtraction} of the gradient! Addition would be gradient ascent (for maximization).
	\end{frame}
	
	\begin{frame}{Q18: Learning Rate - Definition}
		\fontsize{6.5}{7.5}\selectfont
		\textbf{What does the learning rate $\alpha$ control in gradient descent?}
		
		\begin{enumerate}
			\item The direction of the parameter update.
			\item The step size of each parameter update.
			\item The number of iterations.
			\item The initial parameter values.
		\end{enumerate}
		
		\vspace{0.05cm}
		\textcolor{myblue}{\textbf{Answer: B}}
		\vspace{0.05cm}
		\textcolor{myblue}{\textbf{Explanation:}} $\alpha$ controls step size. Larger $\alpha$ = larger steps (faster but risk overshooting). Smaller $\alpha$ = smaller steps (slower but more stable).
		
		\vspace{0.05cm}
		\textcolor{myred}{\textbf{Why Others Are Wrong:}}
		\begin{itemize}
			\item \textcolor{myred}{A:} Direction is determined by the gradient sign, not the learning rate.
			\item \textcolor{myred}{C:} Number of iterations is a separate hyperparameter (epochs).
			\item \textcolor{myred}{D:} Initial parameter values are set separately before optimization.
		\end{itemize}
		
		\textcolor{myred}{\textbf{Exam Trap:}} Learning rate controls \textcolor{myblue}{step size}, not direction! Direction comes from the gradient. This distinction is crucial for understanding optimization.
	\end{frame}
	
	\begin{frame}{Q19: Learning Rate - Too Large}
		\fontsize{6.5}{7.5}\selectfont
		\textbf{What happens if the learning rate is set too large in gradient descent?}
		
		\begin{enumerate}
			\item The algorithm converges very slowly.
			\item The algorithm may oscillate or diverge, failing to converge.
			\item The algorithm always finds the global minimum.
			\item The algorithm becomes deterministic.
		\end{enumerate}
		
		\vspace{0.05cm}
		\textcolor{myblue}{\textbf{Answer: B}}
		\vspace{0.05cm}
		\textcolor{myblue}{\textbf{Explanation:}} Too large $\alpha$ causes overshooting, oscillation, or divergence. The updates become too large and unstable; loss may increase without bound.
		
		\vspace{0.05cm}
		\textcolor{myred}{\textbf{Why Others Are Wrong:}}
		\begin{itemize}
			\item \textcolor{myred}{A:} This describes a learning rate that is too small (slow convergence).
			\item \textcolor{myred}{C:} Too large learning rate prevents convergence; it certainly doesn't guarantee global minimum.
			\item \textcolor{myred}{D:} The algorithm is already deterministic (given fixed initialization and data).
		\end{itemize}
		
		\textcolor{myred}{\textbf{Exam Trap:}} Too large learning rate $\rightarrow$ \textcolor{myred}{divergence or oscillation}! This is a common failure mode. The opposite (too small) gives slow convergence.
	\end{frame}
	
	\begin{frame}{Q20: Learning Rate - Too Small}
		\fontsize{6.5}{7.5}\selectfont
		\textbf{What happens if the learning rate is set too small in gradient descent?}
		
		\begin{enumerate}
			\item The algorithm converges very quickly.
			\item The algorithm converges very slowly and may get stuck.
			\item The algorithm always finds the global minimum instantly.
			\item The algorithm becomes unstable.
		\end{enumerate}
		
		\vspace{0.05cm}
		\textcolor{myblue}{\textbf{Answer: B}}
		\vspace{0.05cm}
		\textcolor{myblue}{\textbf{Explanation:}} Too small $\alpha$ = tiny steps, very slow convergence. May get stuck in local minima or plateaus because updates are too small to escape.
		
		\vspace{0.05cm}
		\textcolor{myred}{\textbf{Why Others Are Wrong:}}
		\begin{itemize}
			\item \textcolor{myred}{A:} False - too small gives slow convergence, not quick.
			\item \textcolor{myred}{C:} False - no algorithm finds global minimum instantly.
			\item \textcolor{myred}{D:} Too small is stable but slow; instability comes from too large learning rate.
		\end{itemize}
		
		\textcolor{myred}{\textbf{Exam Trap:}} Too small learning rate $\rightarrow$ \textcolor{myred}{slow convergence and getting stuck}! Need to balance between too large (divergence) and too small (slow).
	\end{frame}
	
	\begin{frame}{Q21: Batch Gradient Descent}
		\fontsize{6.5}{7.5}\selectfont
		\textbf{What is Batch Gradient Descent?}
		
		\begin{enumerate}
			\item Gradient descent that updates parameters using a single randomly selected data point.
			\item Gradient descent that updates parameters using the entire training dataset.
			\item Gradient descent that updates parameters using a small random subset of data.
			\item Gradient descent that updates parameters without using any data.
		\end{enumerate}
		
		\vspace{0.05cm}
		\textcolor{myblue}{\textbf{Answer: B}}
		\vspace{0.05cm}
		\textcolor{myblue}{\textbf{Explanation:}} Batch GD computes gradient using the entire dataset for each parameter update. Computationally expensive for large datasets.
		
		\vspace{0.05cm}
		\textcolor{myred}{\textbf{Why Others Are Wrong:}}
		\begin{itemize}
			\item \textcolor{myred}{A:} This describes Stochastic Gradient Descent (SGD) - uses one data point.
			\item \textcolor{myred}{C:} This describes Mini-Batch Gradient Descent - uses a subset.
			\item \textcolor{myred}{D:} This describes random search - not gradient descent.
		\end{itemize}
		
		\textcolor{myred}{\textbf{Exam Trap:}} Batch GD uses \textcolor{myblue}{ALL data} for each update! This is the key differentiator from SGD and mini-batch.
	\end{frame}
	
	\begin{frame}{Q22: Stochastic Gradient Descent}
		\fontsize{6.5}{7.5}\selectfont
		\textbf{What is Stochastic Gradient Descent (SGD)?}
		
		\begin{enumerate}
			\item Gradient descent that updates parameters using the entire dataset.
			\item Gradient descent that updates parameters using a single randomly selected data point per iteration.
			\item Gradient descent that updates parameters using a small random subset of data.
			\item Gradient descent that uses a fixed learning rate schedule.
		\end{enumerate}
		
		\vspace{0.05cm}
		\textcolor{myblue}{\textbf{Answer: B}}
		\vspace{0.05cm}
		\textcolor{myblue}{\textbf{Explanation:}} SGD updates parameters using one randomly selected data point per iteration. Fast and helps escape local minima due to gradient noise.
		
		\vspace{0.05cm}
		\textcolor{myred}{\textbf{Why Others Are Wrong:}}
		\begin{itemize}
			\item \textcolor{myred}{A:} This describes Batch GD (uses all data).
			\item \textcolor{myred}{C:} This describes Mini-Batch GD (uses subset).
			\item \textcolor{myred}{D:} This describes learning rate scheduling - a separate concept.
		\end{itemize}
		
		\textcolor{myred}{\textbf{Exam Trap:}} SGD uses \textcolor{myblue}{ONE data point} per update! This is the key differentiator - not batch, not mini-batch.
	\end{frame}
	
	\begin{frame}{Q23: Mini-Batch Gradient Descent}
		\fontsize{6.5}{7.5}\selectfont
		\textbf{What is Mini-Batch Gradient Descent?}
		
		\begin{enumerate}
			\item Gradient descent that updates parameters using the entire dataset.
			\item Gradient descent that updates parameters using a single data point.
			\item Gradient descent that updates parameters using a small random subset (batch) of data.
			\item Gradient descent that uses a fixed number of iterations.
		\end{enumerate}
		
		\vspace{0.05cm}
		\textcolor{myblue}{\textbf{Answer: C}}
		\vspace{0.05cm}
		\textcolor{myblue}{\textbf{Explanation:}} Mini-Batch GD balances Batch GD (stable) and SGD (fast). Uses small batch (e.g., 32, 64, 128 samples) per update.
		
		\vspace{0.05cm}
		\textcolor{myred}{\textbf{Why Others Are Wrong:}}
		\begin{itemize}
			\item \textcolor{myred}{A:} This is Batch GD (all data).
			\item \textcolor{myred}{B:} This is SGD (one data point).
			\item \textcolor{myred}{D:} This describes a stopping criterion, not the gradient method.
		\end{itemize}
		
		\textcolor{myred}{\textbf{Exam Trap:}} Mini-Batch GD uses a \textcolor{myblue}{small random subset} of data! This is the most common in practice - balances stability and speed.
	\end{frame}
	
	\begin{frame}{Q24: Batch vs SGD - Computational Cost}
		\fontsize{6.5}{7.5}\selectfont
		\textbf{Why is Stochastic Gradient Descent (SGD) often preferred over Batch Gradient Descent for large datasets?}
		
		\begin{enumerate}
			\item SGD always converges faster than Batch GD.
			\item SGD is computationally cheaper per iteration because it uses fewer data points.
			\item SGD always finds the global minimum.
			\item SGD does not require a loss function.
		\end{enumerate}
		
		\vspace{0.05cm}
		\textcolor{myblue}{\textbf{Answer: B}}
		\vspace{0.05cm}
		\textcolor{myblue}{\textbf{Explanation:}} SGD cheaper per iteration because uses fewer data points (one or small batch) to compute gradient. Scalable to large datasets.
		
		\vspace{0.05cm}
		\textcolor{myred}{\textbf{Why Others Are Wrong:}}
		\begin{itemize}
			\item \textcolor{myred}{A:} Not always - SGD converges noisily and may require more iterations.
			\item \textcolor{myred}{C:} False - SGD doesn't guarantee global minimum (no method does).
			\item \textcolor{myred}{D:} False - SGD requires a loss function to compute gradients.
		\end{itemize}
		
		\textcolor{myred}{\textbf{Exam Trap:}} SGD preferred for large datasets due to \textcolor{myblue}{lower computational cost per iteration}, not because it always converges faster.
	\end{frame}
	
	\begin{frame}{Q25: Gradient Descent - Local Minima}
		\fontsize{6.5}{7.5}\selectfont
		\textbf{What is a local minimum in gradient descent?}
		
		\begin{enumerate}
			\item The point where the loss function is at its global minimum.
			\item The point where the gradient is zero but the function value is higher than at other nearby points.
			\item The point where the gradient is infinite.
			\item The point where the loss function is at its maximum.
		\end{enumerate}
		
		\vspace{0.05cm}
		\textcolor{myblue}{\textbf{Answer: B}}
		\vspace{0.05cm}
		\textcolor{myblue}{\textbf{Explanation:}} Local minimum: gradient $\approx 0$, function value lower than neighbors, but not globally lowest. Gradient descent can get stuck here.
		
		\vspace{0.05cm}
		\textcolor{myred}{\textbf{Why Others Are Wrong:}}
		\begin{itemize}
			\item \textcolor{myred}{A:} This describes global minimum (lowest point overall).
			\item \textcolor{myred}{C:} Gradient is zero at minima, not infinite.
			\item \textcolor{myred}{D:} This describes a maximum (gradient zero, higher than neighbors).
		\end{itemize}
		
		\textcolor{myred}{\textbf{Exam Trap:}} Local minima are \textcolor{myred}{traps} where gradient descent can get stuck! They have zero gradient but aren't the global best.
	\end{frame}
	
	\begin{frame}{Q26: Saddle Points}
		\fontsize{6.5}{7.5}\selectfont
		\textbf{What is a saddle point in gradient descent?}
		
		\begin{enumerate}
			\item A point where the gradient is zero and the function curves upward in all directions.
			\item A point where the gradient is zero but the function curves upward in some directions and downward in others.
			\item A point where the gradient is infinite.
			\item A point where the function is at its global minimum.
		\end{enumerate}
		
		\vspace{0.05cm}
		\textcolor{myblue}{\textbf{Answer: B}}
		\vspace{0.05cm}
		\textcolor{myblue}{\textbf{Explanation:}} Saddle point: gradient zero, neither local min nor max. Curves upward in some directions, downward in others. Common in high dimensions.
		
		\vspace{0.05cm}
		\textcolor{myred}{\textbf{Why Others Are Wrong:}}
		\begin{itemize}
			\item \textcolor{myred}{A:} This describes a local minimum (curves upward in all directions).
			\item \textcolor{myred}{C:} Gradient is zero at saddle points, not infinite.
			\item \textcolor{myred}{D:} This describes global minimum - saddle points are not minima.
		\end{itemize}
		
		\textcolor{myred}{\textbf{Exam Trap:}} Saddle points have \textcolor{myblue}{zero gradient} but are neither minima nor maxima! More common than local minima in high dimensions.
	\end{frame}
	
	\begin{frame}{Q27: Gradient Descent - Momentum}
		\fontsize{6.5}{7.5}\selectfont
		\textbf{What is the purpose of momentum in gradient descent?}
		
		\begin{enumerate}
			\item To decrease the learning rate over time.
			\item To accelerate convergence and reduce oscillations by accumulating past gradients.
			\item To increase the batch size.
			\item To randomly initialize parameters.
		\end{enumerate}
		
		\vspace{0.05cm}
		\textcolor{myblue}{\textbf{Answer: B}}
		\vspace{0.05cm}
		\textcolor{myblue}{\textbf{Explanation:}} Momentum adds fraction of previous update: $v_t = \beta v_{t-1} + (1-\beta)\nabla L(\theta)$, $\theta = \theta - \alpha v_t$. Smooths oscillations, helps escape shallow minima.
		
		\vspace{0.05cm}
		\textcolor{myred}{\textbf{Why Others Are Wrong:}}
		\begin{itemize}
			\item \textcolor{myred}{A:} This describes learning rate decay - a separate concept.
			\item \textcolor{myred}{C:} This describes batch size - unrelated to momentum.
			\item \textcolor{myred}{D:} This describes parameter initialization - a different step.
		\end{itemize}
		
		\textcolor{myred}{\textbf{Exam Trap:}} Momentum uses \textcolor{myblue}{past gradients} to smooth updates! This helps accelerate convergence and escape local minima.
	\end{frame}
	
	\begin{frame}{Q28: Adam Optimizer}
		\fontsize{6.5}{7.5}\selectfont
		\textbf{What is the Adam optimizer in gradient descent?}
		
		\begin{enumerate}
			\item A method that uses a fixed learning rate.
			\item An adaptive learning rate optimizer that combines momentum and RMSprop.
			\item A method that only uses the first derivative.
			\item A method that randomly initializes parameters.
		\end{enumerate}
		
		\vspace{0.05cm}
		\textcolor{myblue}{\textbf{Answer: B}}
		\vspace{0.05cm}
		\textcolor{myblue}{\textbf{Explanation:}} Adam = Adaptive Moment Estimation. Combines momentum (first moment) and RMSprop (second moment). Adapts learning rate per parameter.
		
		\vspace{0.05cm}
		\textcolor{myred}{\textbf{Why Others Are Wrong:}}
		\begin{itemize}
			\item \textcolor{myred}{A:} Adam adapts learning rate dynamically; it's not fixed.
			\item \textcolor{myred}{C:} Adam uses first AND second moments - more than just first derivative.
			\item \textcolor{myred}{D:} Adam is an optimizer, not an initialization method.
		\end{itemize}
		
		\textcolor{myred}{\textbf{Exam Trap:}} Adam is \textcolor{myblue}{adaptive} - adjusts learning rates per parameter! Most popular optimizer in deep learning.
	\end{frame}
	
	\begin{frame}{Q29: Learning Rate Scheduling}
		\fontsize{6.5}{7.5}\selectfont
		\textbf{What is learning rate scheduling?}
		
		\begin{enumerate}
			\item A method that randomly changes the learning rate.
			\item A method that adjusts the learning rate during training according to a predefined schedule.
			\item A method that keeps the learning rate constant throughout training.
			\item A method that changes the batch size during training.
		\end{enumerate}
		
		\vspace{0.05cm}
		\textcolor{myblue}{\textbf{Answer: B}}
		\vspace{0.05cm}
		\textcolor{myblue}{\textbf{Explanation:}} Learning rate scheduling adjusts LR during training (step decay, exponential decay, cyclical). Larger steps early, smaller steps later.
		
		\vspace{0.05cm}
		\textcolor{myred}{\textbf{Why Others Are Wrong:}}
		\begin{itemize}
			\item \textcolor{myred}{A:} Random changes would be unpredictable and not a "schedule".
			\item \textcolor{myred}{C:} This is a constant learning rate - the opposite of scheduling.
			\item \textcolor{myred}{D:} This describes batch size scheduling - a different concept.
		\end{itemize}
		
		\textcolor{myred}{\textbf{Exam Trap:}} Learning rate schedules \textcolor{myblue}{adjust the learning rate over time}! Not random, not constant, not batch size.
	\end{frame}
	
	\begin{frame}{Q30: Step Decay}
		\fontsize{6.5}{7.5}\selectfont
		\textbf{What is step decay in learning rate scheduling?}
		
		\begin{enumerate}
			\item The learning rate decreases continuously over time.
			\item The learning rate decreases by a factor at specific intervals (e.g., every 10 epochs).
			\item The learning rate increases over time.
			\item The learning rate is randomly changed.
		\end{enumerate}
		
		\vspace{0.05cm}
		\textcolor{myblue}{\textbf{Answer: B}}
		\vspace{0.05cm}
		\textcolor{myblue}{\textbf{Explanation:}} Step decay reduces LR by factor (e.g., multiply by 0.1) at specific intervals (e.g., every 10 epochs or after certain iterations).
		
		\vspace{0.05cm}
		\textcolor{myred}{\textbf{Why Others Are Wrong:}}
		\begin{itemize}
			\item \textcolor{myred}{A:} This describes exponential decay (continuous decrease).
			\item \textcolor{myred}{C:} This describes increasing learning rate (rarely used).
			\item \textcolor{myred}{D:} This describes random changes - not a defined schedule.
		\end{itemize}
		
		\textcolor{myred}{\textbf{Exam Trap:}} Step decay reduces learning rate \textcolor{myblue}{at specific intervals}! This is a key distinction from exponential decay (continuous).
	\end{frame}
	
	\begin{frame}{Q31: Exponential Decay}
		\fontsize{6.5}{7.5}\selectfont
		\textbf{What is exponential decay in learning rate scheduling?}
		
		\begin{enumerate}
			\item $\alpha_t = \alpha_0 \times \gamma^{t}$
			\item $\alpha_t = \alpha_0 / t$
			\item $\alpha_t = \alpha_0 \times (1 - \gamma t)$
			\item $\alpha_t = \alpha_0$ (constant)
		\end{enumerate}
		
		\vspace{0.05cm}
		\textcolor{myblue}{\textbf{Answer: A}}
		\vspace{0.05cm}
		\textcolor{myblue}{\textbf{Explanation:}} Exponential decay: $\alpha_t = \alpha_0 \times \gamma^{t}$ where $\gamma$ is decay rate (typically 0.9-0.99). Smooth continuous decay.
		
		\vspace{0.05cm}
		\textcolor{myred}{\textbf{Why Others Are Wrong:}}
		\begin{itemize}
			\item \textcolor{myred}{B:} This is inverse time decay ($\alpha_t = \alpha_0 / t$) - different schedule.
			\item \textcolor{myred}{C:} This is linear decay - decreases linearly, not exponentially.
			\item \textcolor{myred}{D:} This is constant learning rate - no decay at all.
		\end{itemize}
		
		\textcolor{myred}{\textbf{Exam Trap:}} Exponential decay: $\alpha_t = \alpha_0 \times \gamma^{t}$! Know the formula and how it differs from step decay and inverse decay.
	\end{frame}
	
	\begin{frame}{Q32: Gradient Descent - Convergence}
		\fontsize{6.5}{7.5}\selectfont
		\textbf{What does convergence mean in gradient descent?}
		
		\begin{enumerate}
			\item The algorithm has found the global minimum.
			\item The algorithm has stopped updating parameters.
			\item The loss function has stopped decreasing significantly (or the gradient is near zero).
			\item The algorithm has reached the maximum number of iterations.
		\end{enumerate}
		
		\vspace{0.05cm}
		\textcolor{myblue}{\textbf{Answer: C}}
		\vspace{0.05cm}
		\textcolor{myblue}{\textbf{Explanation:}} Convergence = loss stopped decreasing significantly or gradient near zero. Found a minimum (local or global).
		
		\vspace{0.05cm}
		\textcolor{myred}{\textbf{Why Others Are Wrong:}}
		\begin{itemize}
			\item \textcolor{myred}{A:} Ideal but not guaranteed - could be local minimum.
			\item \textcolor{myred}{B:} Could be due to learning rate = 0, not convergence.
			\item \textcolor{myred}{D:} This is a stopping criterion, not convergence.
		\end{itemize}
		
		\textcolor{myred}{\textbf{Exam Trap:}} Convergence means \textcolor{myblue}{loss has stopped decreasing} significantly! Not necessarily global minimum, not just stopping updates.
	\end{frame}
	
	\begin{frame}{Q33: Gradient Descent - Loss Landscape}
		\fontsize{6.5}{7.5}\selectfont
		\textbf{In high-dimensional spaces (e.g., neural networks), which is more common: local minima or saddle points?}
		
		\begin{enumerate}
			\item Local minima are more common.
			\item Saddle points are more common.
			\item Global minima are always found.
			\item Neither local minima nor saddle points exist.
		\end{enumerate}
		
		\vspace{0.05cm}
		\textcolor{myblue}{\textbf{Answer: B}}
		\vspace{0.05cm}
		\textcolor{myblue}{\textbf{Explanation:}} In high dimensions, saddle points are more common than local minima. Research shows local minima in neural network landscapes are relatively rare.
		
		\vspace{0.05cm}
		\textcolor{myred}{\textbf{Why Others Are Wrong:}}
		\begin{itemize}
			\item \textcolor{myred}{A:} False - local minima are rare in high dimensions.
			\item \textcolor{myred}{C:} False - no guarantee of finding global minimum.
			\item \textcolor{myred}{D:} False - both exist, but saddle points are more common.
		\end{itemize}
		
		\textcolor{myred}{\textbf{Exam Trap:}} Saddle points are \textcolor{myblue}{more common} than local minima in high dimensions! This is a key insight for deep learning optimization.
	\end{frame}
	
	\begin{frame}{Q34: Gradient Descent - Vanishing Gradients}
		\fontsize{6.5}{7.5}\selectfont
		\textbf{What is the vanishing gradient problem in deep neural networks?}
		
		\begin{enumerate}
			\item Gradients become very large during backpropagation.
			\item Gradients become very small during backpropagation, preventing early layers from learning.
			\item Gradients oscillate wildly during training.
			\item Gradients are zero for all parameters.
		\end{enumerate}
		
		\vspace{0.05cm}
		\textcolor{myblue}{\textbf{Answer: B}}
		\vspace{0.05cm}
		\textcolor{myblue}{\textbf{Explanation:}} Vanishing gradients: gradients become very small during backpropagation, especially in early layers. Prevents effective updates. Often caused by sigmoid/tanh in saturation regions.
		
		\vspace{0.05cm}
		\textcolor{myred}{\textbf{Why Others Are Wrong:}}
		\begin{itemize}
			\item \textcolor{myred}{A:} This describes exploding gradients (large gradients).
			\item \textcolor{myred}{C:} This describes oscillations - a different problem.
			\item \textcolor{myred}{D:} Extreme case - gradients becoming exactly zero everywhere.
		\end{itemize}
		
		\textcolor{myred}{\textbf{Exam Trap:}} Vanishing gradients = \textcolor{myblue}{very small gradients} in early layers! Use ReLU activation to mitigate.
	\end{frame}
	
	\begin{frame}{Q35: Gradient Descent - Exploding Gradients}
		\fontsize{6.5}{7.5}\selectfont
		\textbf{What is the exploding gradient problem in deep neural networks?}
		
		\begin{enumerate}
			\item Gradients become very small during backpropagation.
			\item Gradients become very large during backpropagation, causing unstable updates.
			\item Gradients are zero for all parameters.
			\item Gradients oscillate in a regular pattern.
		\end{enumerate}
		
		\vspace{0.05cm}
		\textcolor{myblue}{\textbf{Answer: B}}
		\vspace{0.05cm}
		\textcolor{myblue}{\textbf{Explanation:}} Exploding gradients: gradients become very large during backpropagation, leading to large parameter updates and unstable training. Loss may diverge.
		
		\vspace{0.05cm}
		\textcolor{myred}{\textbf{Why Others Are Wrong:}}
		\begin{itemize}
			\item \textcolor{myred}{A:} This describes vanishing gradients (small gradients).
			\item \textcolor{myred}{C:} This is the extreme opposite - zero gradients.
			\item \textcolor{myred}{D:} Regular oscillations are a different phenomenon.
		\end{itemize}
		
		\textcolor{myred}{\textbf{Exam Trap:}} Exploding gradients = \textcolor{myblue}{very large gradients} causing instability! Use gradient clipping to mitigate.
	\end{frame}
	
	% ============================================================
	% SECTION 3: BACKPROPAGATION - QUESTIONS 36-45
	% ============================================================
	
	\begin{frame}{Q36: Backpropagation - Definition}
		\fontsize{6.5}{7.5}\selectfont
		\textbf{What is backpropagation in neural networks?}
		
		\begin{enumerate}
			\item A method to initialize neural network weights.
			\item An algorithm that computes gradients of the loss function with respect to all weights using the chain rule.
			\item A method to choose the number of hidden layers.
			\item A method to select the activation function.
		\end{enumerate}
		
		\vspace{0.05cm}
		\textcolor{myblue}{\textbf{Answer: B}}
		\vspace{0.05cm}
		\textcolor{myblue}{\textbf{Explanation:}} Backpropagation computes gradients of loss wrt all weights using chain rule. Propagates errors backward from output to input layers.
		
		\vspace{0.05cm}
		\textcolor{myred}{\textbf{Why Others Are Wrong:}}
		\begin{itemize}
			\item \textcolor{myred}{A:} This describes weight initialization - a separate step.
			\item \textcolor{myred}{C:} This describes architecture selection - a different problem.
			\item \textcolor{myred}{D:} This describes activation function selection - not backpropagation.
		\end{itemize}
		
		\textcolor{myred}{\textbf{Exam Trap:}} Backpropagation uses the \textcolor{myblue}{chain rule} to compute gradients! This is the fundamental algorithm for neural network training.
	\end{frame}
	
	\begin{frame}{Q37: Backpropagation - Chain Rule}
		\fontsize{6.5}{7.5}\selectfont
		\textbf{Which mathematical rule is the foundation of backpropagation?}
		
		\begin{enumerate}
			\item The product rule
			\item The quotient rule
			\item The chain rule
			\item The power rule
		\end{enumerate}
		
		\vspace{0.05cm}
		\textcolor{myblue}{\textbf{Answer: C}}
		\vspace{0.05cm}
		\textcolor{myblue}{\textbf{Explanation:}} Backpropagation is based on chain rule of calculus. Loss function is composition of many functions (layers, activations), chain rule propagates gradients backward.
		
		\vspace{0.05cm}
		\textcolor{myred}{\textbf{Why Others Are Wrong:}}
		\begin{itemize}
			\item \textcolor{myred}{A:} Product rule applies to products of functions - not the core of backprop.
			\item \textcolor{myred}{B:} Quotient rule applies to ratios - not the core of backprop.
			\item \textcolor{myred}{D:} Power rule applies to powers - not the core of backprop.
		\end{itemize}
		
		\textcolor{myred}{\textbf{Exam Trap:}} Backpropagation = \textcolor{myblue}{chain rule} applied to neural networks! This is the single most important mathematical concept for neural network training.
	\end{frame}
	
	\begin{frame}{Q38: Backpropagation - Forward Pass}
		\fontsize{6.5}{7.5}\selectfont
		\textbf{What happens during the forward pass in backpropagation?}
		
		\begin{enumerate}
			\item Gradients are computed and weights are updated.
			\item Input data is passed through the network to compute the output and loss.
			\item The loss function is minimized.
			\item The weights are initialized.
		\end{enumerate}
		
		\vspace{0.05cm}
		\textcolor{myblue}{\textbf{Answer: B}}
		\vspace{0.05cm}
		\textcolor{myblue}{\textbf{Explanation:}} Forward pass: input $\to$ hidden layers $\to$ output, compute output and loss. First step in backpropagation.
		
		\vspace{0.05cm}
		\textcolor{myred}{\textbf{Why Others Are Wrong:}}
		\begin{itemize}
			\item \textcolor{myred}{A:} This describes backward pass (gradient computation and update).
			\item \textcolor{myred}{C:} Optimization happens over many iterations, not just the forward pass.
			\item \textcolor{myred}{D:} Initialization happens before training begins.
		\end{itemize}
		
		\textcolor{myred}{\textbf{Exam Trap:}} Forward pass computes \textcolor{myblue}{output and loss}; backward pass computes gradients! The order is critical.
	\end{frame}
	
	\begin{frame}{Q39: Backpropagation - Backward Pass}
		\fontsize{6.5}{7.5}\selectfont
		\textbf{What happens during the backward pass in backpropagation?}
		
		\begin{enumerate}
			\item Input data is passed through the network.
			\item Gradients of the loss with respect to each weight are computed using the chain rule.
			\item The learning rate is adjusted.
			\item The loss function is evaluated.
		\end{enumerate}
		
		\vspace{0.05cm}
		\textcolor{myblue}{\textbf{Answer: B}}
		\vspace{0.05cm}
		\textcolor{myblue}{\textbf{Explanation:}} Backward pass computes gradients of loss wrt each weight using chain rule. Propagates errors from output layer back to input layer.
		
		\vspace{0.05cm}
		\textcolor{myred}{\textbf{Why Others Are Wrong:}}
		\begin{itemize}
			\item \textcolor{myred}{A:} This describes forward pass (input through network).
			\item \textcolor{myred}{C:} Learning rate adjustment is a separate process (scheduling).
			\item \textcolor{myred}{D:} Loss evaluation happens during forward pass.
		\end{itemize}
		
		\textcolor{myred}{\textbf{Exam Trap:}} Backward pass computes \textcolor{myblue}{gradients} of all weights! The name "backward" comes from propagating errors backward from output to input.
	\end{frame}
	
	\begin{frame}{Q40: Backpropagation - Weight Update}
		\fontsize{6.5}{7.5}\selectfont
		\textbf{After computing gradients in backpropagation, how are the weights updated?}
		
		\begin{enumerate}
			\item $w_{new} = w_{old} + \alpha \nabla L(w)$
			\item $w_{new} = w_{old} - \alpha \nabla L(w)$
			\item $w_{new} = \nabla L(w)$
			\item $w_{new} = w_{old} \times \alpha$
		\end{enumerate}
		
		\vspace{0.05cm}
		\textcolor{myblue}{\textbf{Answer: B}}
		\vspace{0.05cm}
		\textcolor{myblue}{\textbf{Explanation:}} $w_{new} = w_{old} - \alpha \nabla L(w)$. Subtract gradient to move in direction of decreasing loss.
		
		\vspace{0.05cm}
		\textcolor{myred}{\textbf{Why Others Are Wrong:}}
		\begin{itemize}
			\item \textcolor{myred}{A:} This is gradient ascent - would increase loss (maximization).
			\item \textcolor{myred}{C:} Ignores learning rate and old weights - not a valid update.
			\item \textcolor{myred}{D:} Multiplication alone doesn't use gradient information.
		\end{itemize}
		
		\textcolor{myred}{\textbf{Exam Trap:}} Weights updated by \textcolor{myblue}{subtracting} gradient $\times$ learning rate! This is the gradient descent update rule.
	\end{frame}
	
	\begin{frame}{Q41: Backpropagation - Computational Graph}
		\fontsize{6.5}{7.5}\selectfont
		\textbf{What is a computational graph in the context of backpropagation?}
		
		\begin{enumerate}
			\item A graph showing the architecture of the neural network.
			\item A graph representing the forward computation as a directed acyclic graph (DAG) that enables efficient gradient computation.
			\item A graph showing the training data distribution.
			\item A graph showing the loss function.
		\end{enumerate}
		
		\vspace{0.05cm}
		\textcolor{myblue}{\textbf{Answer: B}}
		\vspace{0.05cm}
		\textcolor{myblue}{\textbf{Explanation:}} Computational graph represents forward computation as DAG of operations. Enables efficient gradient computation via automatic differentiation.
		
		\vspace{0.05cm}
		\textcolor{myred}{\textbf{Why Others Are Wrong:}}
		\begin{itemize}
			\item \textcolor{myred}{A:} This is a network architecture diagram - not the computational graph.
			\item \textcolor{myred}{C:} This is data distribution - unrelated to computational graph.
			\item \textcolor{myred}{D:} This is a loss function plot - not the computational graph.
		\end{itemize}
		
		\textcolor{myred}{\textbf{Exam Trap:}} Computational graphs enable \textcolor{myblue}{automatic differentiation} and backpropagation! Foundation of deep learning frameworks like TensorFlow and PyTorch.
	\end{frame}
	
	\begin{frame}{Q42: Backpropagation - Auto Differentiation}
		\fontsize{6.5}{7.5}\selectfont
		\textbf{What is automatic differentiation in deep learning frameworks?}
		
		\begin{enumerate}
			\item Automatic selection of the learning rate.
			\item Automatic computation of gradients using the chain rule on computational graphs.
			\item Automatic selection of the number of layers.
			\item Automatic data preprocessing.
		\end{enumerate}
		
		\vspace{0.05cm}
		\textcolor{myblue}{\textbf{Answer: B}}
		\vspace{0.05cm}
		\textcolor{myblue}{\textbf{Explanation:}} Automatic differentiation (autodiff) = automatic gradient computation using chain rule on computational graphs. Frameworks use autodiff to implement backpropagation.
		
		\vspace{0.05cm}
		\textcolor{myred}{\textbf{Why Others Are Wrong:}}
		\begin{itemize}
			\item \textcolor{myred}{A:} Learning rate selection is hyperparameter tuning - separate concept.
			\item \textcolor{myred}{C:} Architecture selection is a separate problem (NAS).
			\item \textcolor{myred}{D:} Data preprocessing happens before model training.
		\end{itemize}
		
		\textcolor{myred}{\textbf{Exam Trap:}} Automatic differentiation is how frameworks implement \textcolor{myblue}{backpropagation}! This is the key technical insight.
	\end{frame}
	
	\begin{frame}{Q43: Backpropagation - Memory Requirements}
		\fontsize{6.5}{7.5}\selectfont
		\textbf{What is a key memory requirement for backpropagation?}
		
		\begin{enumerate}
			\item Backpropagation requires storing all intermediate activations from the forward pass.
			\item Backpropagation requires no memory beyond the model parameters.
			\item Backpropagation requires storing only the final loss.
			\item Backpropagation requires storing only the gradients.
		\end{enumerate}
		
		\vspace{0.05cm}
		\textcolor{myblue}{\textbf{Answer: A}}
		\vspace{0.05cm}
		\textcolor{myblue}{\textbf{Explanation:}} Backpropagation requires storing all intermediate activations (outputs of each layer) from forward pass. Needed during backward pass to compute gradients.
		
		\vspace{0.05cm}
		\textcolor{myred}{\textbf{Why Others Are Wrong:}}
		\begin{itemize}
			\item \textcolor{myred}{B:} False - memory is needed for intermediate activations.
			\item \textcolor{myred}{C:} False - final loss alone is insufficient for gradient computation.
			\item \textcolor{myred}{D:} False - gradients are computed during backward pass, not stored from forward pass.
		\end{itemize}
		
		\textcolor{myred}{\textbf{Exam Trap:}} Backpropagation requires storing \textcolor{myblue}{all intermediate activations}! This is why deep networks are memory-intensive.
	\end{frame}
	
	\begin{frame}{Q44: Backpropagation - Training Loop}
		\fontsize{6.5}{7.5}\selectfont
		\textbf{What is the correct order of operations in a training loop with backpropagation?}
		
		\begin{enumerate}
			\item Forward pass $\to$ Compute loss $\to$ Backward pass $\to$ Update weights
			\item Backward pass $\to$ Compute loss $\to$ Forward pass $\to$ Update weights
			\item Compute loss $\to$ Forward pass $\to$ Backward pass $\to$ Update weights
			\item Update weights $\to$ Forward pass $\to$ Compute loss $\to$ Backward pass
		\end{enumerate}
		
		\vspace{0.05cm}
		\textcolor{myblue}{\textbf{Answer: A}}
		\vspace{0.05cm}
		\textcolor{myblue}{\textbf{Explanation:}} Correct order: (1) Forward pass: compute predictions and loss, (2) Compute loss, (3) Backward pass: compute gradients, (4) Update weights.
		
		\vspace{0.05cm}
		\textcolor{myred}{\textbf{Why Others Are Wrong:}}
		\begin{itemize}
			\item \textcolor{myred}{B:} Backward pass cannot happen before forward pass and loss computation.
			\item \textcolor{myred}{C:} Loss cannot be computed before forward pass.
			\item \textcolor{myred}{D:} Weights cannot be updated before computing gradients.
		\end{itemize}
		
		\textcolor{myred}{\textbf{Exam Trap:}} Forward $\to$ Loss $\to$ Backward $\to$ Update! This order is absolutely critical and tested frequently.
	\end{frame}
	
	\begin{frame}{Q45: Backpropagation - Summary}
		\fontsize{6.5}{7.5}\selectfont
		\textbf{Which of the following statements about backpropagation is FALSE?}
		
		\begin{enumerate}
			\item Backpropagation computes gradients of the loss with respect to all weights.
			\item Backpropagation uses the chain rule of calculus.
			\item Backpropagation requires storing intermediate activations.
			\item Backpropagation always finds the global minimum of the loss function.
		\end{enumerate}
		
		\vspace{0.05cm}
		\textcolor{myblue}{\textbf{Answer: D}}
		\vspace{0.05cm}
		\textcolor{myblue}{\textbf{Explanation:}} Backpropagation does NOT guarantee finding the global minimum. Can get stuck in local minima or saddle points.
		
		\vspace{0.05cm}
		\textcolor{myred}{\textbf{Why Others Are Wrong:}}
		\begin{itemize}
			\item \textcolor{myred}{A:} True - this is exactly what backpropagation does.
			\item \textcolor{myred}{B:} True - chain rule is the mathematical foundation.
			\item \textcolor{myred}{C:} True - intermediate activations are stored for gradient computation.
		\end{itemize}
		
		\textcolor{myred}{\textbf{Exam Trap:}} Backpropagation does \textcolor{myred}{NOT} guarantee finding the global minimum! This is a common misconception.
	\end{frame}
	
	% ============================================================
	% SECTION 4: BIAS-VARIANCE TRADEOFF - QUESTIONS 46-60
	% ============================================================
	
	\begin{frame}{Q46: Bias-Variance Tradeoff - Definition}
		\fontsize{6.5}{7.5}\selectfont
		\textbf{What is the bias-variance tradeoff in supervised learning?}
		
		\begin{enumerate}
			\item The tradeoff between model complexity and interpretability.
			\item The tradeoff between the model's ability to fit the training data (bias) and its ability to generalize to new data (variance).
			\item The tradeoff between training speed and accuracy.
			\item The tradeoff between the number of features and the number of observations.
		\end{enumerate}
		
		\vspace{0.05cm}
		\textcolor{myblue}{\textbf{Answer: B}}
		\vspace{0.05cm}
		\textcolor{myblue}{\textbf{Explanation:}} Bias-variance tradeoff: tension between fitting training data (low bias) and generalizing to new data (low variance). Simple models: high bias, low variance. Complex models: low bias, high variance.
		
		\vspace{0.05cm}
		\textcolor{myred}{\textbf{Why Others Are Wrong:}}
		\begin{itemize}
			\item \textcolor{myred}{A:} This is the accuracy-interpretability tradeoff - a different concept.
			\item \textcolor{myred}{C:} This is the speed-accuracy tradeoff - not bias-variance.
			\item \textcolor{myred}{D:} This describes the curse of dimensionality - a different concept.
		\end{itemize}
		
		\textcolor{myred}{\textbf{Exam Trap:}} Bias = \textcolor{myblue}{error from assumptions}; Variance = \textcolor{myblue}{error from sensitivity to training data}! Understanding this distinction is critical.
	\end{frame}
	
	\begin{frame}{Q47: Bias - Definition}
		\fontsize{6.5}{7.5}\selectfont
		\textbf{What does "bias" refer to in the bias-variance tradeoff?}
		
		\begin{enumerate}
			\item The error introduced by the model's assumptions and limitations.
			\item The error introduced by using a limited training sample.
			\item The error introduced by noise in the data.
			\item The error introduced by the optimization algorithm.
		\end{enumerate}
		
		\vspace{0.05cm}
		\textcolor{myblue}{\textbf{Answer: A}}
		\vspace{0.05cm}
		\textcolor{myblue}{\textbf{Explanation:}} Bias = error from model's assumptions and limitations. High bias = strong assumptions, underfitting (e.g., linear regression for nonlinear data).
		
		\vspace{0.05cm}
		\textcolor{myred}{\textbf{Why Others Are Wrong:}}
		\begin{itemize}
			\item \textcolor{myred}{B:} This describes variance (sensitivity to training sample).
			\item \textcolor{myred}{C:} This describes irreducible error (noise in data).
			\item \textcolor{myred}{D:} This describes optimization error - a different concept.
		\end{itemize}
		
		\textcolor{myred}{\textbf{Exam Trap:}} Bias = \textcolor{myblue}{error from model assumptions} (underfitting)! Not from data noise or sampling.
	\end{frame}
	
	\begin{frame}{Q48: Variance - Definition}
		\fontsize{6.5}{7.5}\selectfont
		\textbf{What does "variance" refer to in the bias-variance tradeoff?}
		
		\begin{enumerate}
			\item The error introduced by the model's assumptions.
			\item The error introduced by the model's sensitivity to fluctuations in the training data.
			\item The error introduced by noise in the data.
			\item The error introduced by the choice of loss function.
		\end{enumerate}
		
		\vspace{0.05cm}
		\textcolor{myblue}{\textbf{Answer: B}}
		\vspace{0.05cm}
		\textcolor{myblue}{\textbf{Explanation:}} Variance = error from model's sensitivity to training data fluctuations. High variance = overfitting (fits training data too closely, fails on new data).
		
		\vspace{0.05cm}
		\textcolor{myred}{\textbf{Why Others Are Wrong:}}
		\begin{itemize}
			\item \textcolor{myred}{A:} This describes bias (model assumptions).
			\item \textcolor{myred}{C:} This describes irreducible error (data noise).
			\item \textcolor{myred}{D:} This describes a different type of error.
		\end{itemize}
		
		\textcolor{myred}{\textbf{Exam Trap:}} Variance = \textcolor{myblue}{error from sensitivity to training data} (overfitting)! Not from model assumptions or data noise.
	\end{frame}
	
	\begin{frame}{Q49: High Bias - Characteristics}
		\fontsize{6.5}{7.5}\selectfont
		\textbf{Which of the following is a characteristic of a model with high bias?}
		
		\begin{enumerate}
			\item It fits the training data very well.
			\item It performs poorly on both training and test data (underfitting).
			\item It performs well on training data but poorly on test data.
			\item It has high variance.
		\end{enumerate}
		
		\vspace{0.05cm}
		\textcolor{myblue}{\textbf{Answer: B}}
		\vspace{0.05cm}
		\textcolor{myblue}{\textbf{Explanation:}} High bias (underfitting): performs poorly on both training and test data. Model has made overly strong assumptions and failed to capture patterns.
		
		\vspace{0.05cm}
		\textcolor{myred}{\textbf{Why Others Are Wrong:}}
		\begin{itemize}
			\item \textcolor{myred}{A:} This describes low bias (fits training data well).
			\item \textcolor{myred}{C:} This describes high variance (overfitting).
			\item \textcolor{myred}{D:} High bias is associated with low variance, not high.
		\end{itemize}
		
		\textcolor{myred}{\textbf{Exam Trap:}} High bias = \textcolor{myblue}{poor performance on both training and test} (underfitting)! Not just poor test performance.
	\end{frame}
	
	\begin{frame}{Q50: High Variance - Characteristics}
		\fontsize{6.5}{7.5}\selectfont
		\textbf{Which of the following is a characteristic of a model with high variance?}
		
		\begin{enumerate}
			\item It performs poorly on both training and test data.
			\item It performs very well on training data but poorly on test data (overfitting).
			\item It has low complexity.
			\item It has high bias.
		\end{enumerate}
		
		\vspace{0.05cm}
		\textcolor{myblue}{\textbf{Answer: B}}
		\vspace{0.05cm}
		\textcolor{myblue}{\textbf{Explanation:}} High variance (overfitting): performs very well on training data (memorizes patterns) but poorly on test data (fails to generalize).
		
		\vspace{0.05cm}
		\textcolor{myred}{\textbf{Why Others Are Wrong:}}
		\begin{itemize}
			\item \textcolor{myred}{A:} This describes high bias (underfitting).
			\item \textcolor{myred}{C:} High variance is associated with high complexity, not low.
			\item \textcolor{myred}{D:} High variance is associated with low bias, not high.
		\end{itemize}
		
		\textcolor{myred}{\textbf{Exam Trap:}} High variance = \textcolor{myblue}{good training, poor test} (overfitting)! This is the classic signature of overfitting.
	\end{frame}
	
	\begin{frame}{Q51: Underfitting - Definition}
		\fontsize{6.5}{7.5}\selectfont
		\textbf{What is underfitting in supervised learning?}
		
		\begin{enumerate}
			\item When the model is too complex and learns noise in the data.
			\item When the model is too simple to capture the underlying patterns in the data.
			\item When the model performs well on training data but poorly on test data.
			\item When the model has too many parameters.
		\end{enumerate}
		
		\vspace{0.05cm}
		\textcolor{myblue}{\textbf{Answer: B}}
		\vspace{0.05cm}
		\textcolor{myblue}{\textbf{Explanation:}} Underfitting: model too simple (e.g., linear for nonlinear data). Fails to capture underlying patterns. Poor performance on both training and test data.
		
		\vspace{0.05cm}
		\textcolor{myred}{\textbf{Why Others Are Wrong:}}
		\begin{itemize}
			\item \textcolor{myred}{A:} This describes overfitting (too complex, learns noise).
			\item \textcolor{myred}{C:} This describes overfitting (good training, poor test).
			\item \textcolor{myred}{D:} Too many parameters leads to overfitting, not underfitting.
		\end{itemize}
		
		\textcolor{myred}{\textbf{Exam Trap:}} Underfitting = \textcolor{myblue}{model too simple} = poor training and test performance! Not just poor test performance.
	\end{frame}
	
	\begin{frame}{Q52: Overfitting - Definition}
		\fontsize{6.5}{7.5}\selectfont
		\textbf{What is overfitting in supervised learning?}
		
		\begin{enumerate}
			\item When the model is too simple to capture the underlying patterns.
			\item When the model is too complex and learns the noise in the training data.
			\item When the model performs poorly on both training and test data.
			\item When the model has too few parameters.
		\end{enumerate}
		
		\vspace{0.05cm}
		\textcolor{myblue}{\textbf{Answer: B}}
		\vspace{0.05cm}
		\textcolor{myblue}{\textbf{Explanation:}} Overfitting: model too complex (e.g., high-degree polynomial). Learns noise in training data. High training accuracy but poor test accuracy.
		
		\vspace{0.05cm}
		\textcolor{myred}{\textbf{Why Others Are Wrong:}}
		\begin{itemize}
			\item \textcolor{myred}{A:} This describes underfitting (too simple).
			\item \textcolor{myred}{C:} This describes underfitting (poor both).
			\item \textcolor{myred}{D:} Too few parameters leads to underfitting, not overfitting.
		\end{itemize}
		
		\textcolor{myred}{\textbf{Exam Trap:}} Overfitting = \textcolor{myblue}{model too complex} = good training, poor test! This is the classic signature.
	\end{frame}
	
	\begin{frame}{Q53: Causes of Overfitting}
		\fontsize{6.5}{7.5}\selectfont
		\textbf{Which of the following is a common cause of overfitting?}
		
		\begin{enumerate}
			\item Too few features.
			\item Too many parameters relative to the number of training samples.
			\item Too much regularization.
			\item Too simple a model.
		\end{enumerate}
		
		\vspace{0.05cm}
		\textcolor{myblue}{\textbf{Answer: B}}
		\vspace{0.05cm}
		\textcolor{myblue}{\textbf{Explanation:}} Overfitting caused by too many parameters (or too complex model) relative to training samples. Model memorizes training data instead of learning general patterns.
		
		\vspace{0.05cm}
		\textcolor{myred}{\textbf{Why Others Are Wrong:}}
		\begin{itemize}
			\item \textcolor{myred}{A:} Too few features leads to underfitting, not overfitting.
			\item \textcolor{myred}{C:} Too much regularization leads to underfitting.
			\item \textcolor{myred}{D:} Too simple a model leads to underfitting.
		\end{itemize}
		
		\textcolor{myred}{\textbf{Exam Trap:}} Overfitting happens when \textcolor{myblue}{too many parameters for too few samples}! This is the key cause.
	\end{frame}
	
	\begin{frame}{Q54: Causes of Underfitting}
		\fontsize{6.5}{7.5}\selectfont
		\textbf{Which of the following is a common cause of underfitting?}
		
		\begin{enumerate}
			\item Too many parameters.
			\item Insufficient model complexity (e.g., linear model for nonlinear data).
			\item Too little regularization.
			\item Too much training data.
		\end{enumerate}
		
		\vspace{0.05cm}
		\textcolor{myblue}{\textbf{Answer: B}}
		\vspace{0.05cm}
		\textcolor{myblue}{\textbf{Explanation:}} Underfitting caused by insufficient model complexity. Example: linear model for nonlinear data. Model can't capture patterns.
		
		\vspace{0.05cm}
		\textcolor{myred}{\textbf{Why Others Are Wrong:}}
		\begin{itemize}
			\item \textcolor{myred}{A:} Too many parameters leads to overfitting.
			\item \textcolor{myred}{C:} Too little regularization leads to overfitting.
			\item \textcolor{myred}{D:} More training data generally helps, not causes underfitting.
		\end{itemize}
		
		\textcolor{myred}{\textbf{Exam Trap:}} Underfitting happens when \textcolor{myblue}{model is too simple} for the data! This is the key cause.
	\end{frame}
	
	\begin{frame}{Q55: Bias-Variance Tradeoff - Model Complexity}
		\fontsize{6.5}{7.5}\selectfont
		\textbf{As model complexity increases, what happens to bias and variance?}
		
		\begin{enumerate}
			\item Bias decreases and variance decreases.
			\item Bias decreases and variance increases.
			\item Bias increases and variance decreases.
			\item Bias increases and variance increases.
		\end{enumerate}
		
		\vspace{0.05cm}
		\textcolor{myblue}{\textbf{Answer: B}}
		\vspace{0.05cm}
		\textcolor{myblue}{\textbf{Explanation:}} As complexity increases: bias decreases (captures more complex patterns), variance increases (more sensitive to training data). This is the tradeoff.
		
		\vspace{0.05cm}
		\textcolor{myred}{\textbf{Why Others Are Wrong:}}
		\begin{itemize}
			\item \textcolor{myred}{A:} Both decrease is incorrect - variance increases with complexity.
			\item \textcolor{myred}{C:} This is the opposite direction (complexity increases $\to$ bias $\downarrow$, variance $\uparrow$).
			\item \textcolor{myred}{D:} Both increase is incorrect - bias decreases with complexity.
		\end{itemize}
		
		\textcolor{myred}{\textbf{Exam Trap:}} More complexity $\rightarrow$ \textcolor{myblue}{lower bias, higher variance}! This is the fundamental tradeoff direction.
	\end{frame}
	
	\begin{frame}{Q56: Optimal Model Complexity}
		\fontsize{6.5}{7.5}\selectfont
		\textbf{Where is the optimal model complexity located in the bias-variance tradeoff?}
		
		\begin{enumerate}
			\item At the point of highest variance.
			\item At the point where total error (bias$^2$ + variance + irreducible error) is minimized.
			\item At the point of lowest bias.
			\item At the point of highest bias.
		\end{enumerate}
		
		\vspace{0.05cm}
		\textcolor{myblue}{\textbf{Answer: B}}
		\vspace{0.05cm}
		\textcolor{myblue}{\textbf{Explanation:}} Optimal complexity at minimum total error: bias$^2$ + variance + irreducible error. Balances underfitting and overfitting.
		
		\vspace{0.05cm}
		\textcolor{myred}{\textbf{Why Others Are Wrong:}}
		\begin{itemize}
			\item \textcolor{myred}{A:} Highest variance = overfitting (too complex).
			\item \textcolor{myred}{C:} Lowest bias = overfitting (too complex).
			\item \textcolor{myred}{D:} Highest bias = underfitting (too simple).
		\end{itemize}
		
		\textcolor{myred}{\textbf{Exam Trap:}} Optimal complexity = \textcolor{myblue}{minimum total error} (bias$^2$ + variance)! Not at extremes.
	\end{frame}
	
	\begin{frame}{Q57: Bias-Variance Tradeoff - Example}
		\fontsize{6.5}{7.5}\selectfont
		\textbf{A simple linear regression model for a nonlinear relationship would likely have:}
		
		\begin{enumerate}
			\item High bias, low variance (underfitting).
			\item Low bias, high variance (overfitting).
			\item Low bias, low variance (ideal).
			\item High bias, high variance.
		\end{enumerate}
		
		\vspace{0.05cm}
		\textcolor{myblue}{\textbf{Answer: A}}
		\vspace{0.05cm}
		\textcolor{myblue}{\textbf{Explanation:}} Linear model for nonlinear data is too simple. High bias (can't capture nonlinear patterns), low variance (not sensitive to training data). Underfitting.
		
		\vspace{0.05cm}
		\textcolor{myred}{\textbf{Why Others Are Wrong:}}
		\begin{itemize}
			\item \textcolor{myred}{B:} This describes complex model (e.g., high-degree polynomial).
			\item \textcolor{myred}{C:} Ideal but unrealistic for mismatched model.
			\item \textcolor{myred}{D:} Unlikely combination for linear model on nonlinear data.
		\end{itemize}
		
		\textcolor{myred}{\textbf{Exam Trap:}} Too simple $\rightarrow$ \textcolor{myblue}{high bias, low variance} (underfitting)! This is the signature of underfitting.
	\end{frame}
	
	\begin{frame}{Q58: Bias-Variance Tradeoff - High-Degree Polynomial}
		\fontsize{6.5}{7.5}\selectfont
		\textbf{A high-degree polynomial regression model on a small dataset would likely have:}
		
		\begin{enumerate}
			\item High bias, low variance (underfitting).
			\item Low bias, high variance (overfitting).
			\item Low bias, low variance (ideal).
			\item High bias, high variance.
		\end{enumerate}
		
		\vspace{0.05cm}
		\textcolor{myblue}{\textbf{Answer: B}}
		\vspace{0.05cm}
		\textcolor{myblue}{\textbf{Explanation:}} High-degree polynomial on small dataset is too complex. Low bias (fits training data very well), high variance (very sensitive to training data). Overfitting.
		
		\vspace{0.05cm}
		\textcolor{myred}{\textbf{Why Others Are Wrong:}}
		\begin{itemize}
			\item \textcolor{myred}{A:} This describes simple model (underfitting).
			\item \textcolor{myred}{C:} Ideal but unrealistic for complex model on small dataset.
			\item \textcolor{myred}{D:} Unlikely combination for high-degree polynomial.
		\end{itemize}
		
		\textcolor{myred}{\textbf{Exam Trap:}} Too complex $\rightarrow$ \textcolor{myblue}{low bias, high variance} (overfitting)! This is the signature of overfitting.
	\end{frame}
	
	\begin{frame}{Q59: Bias-Variance Tradeoff - Irreducible Error}
		\fontsize{6.5}{7.5}\selectfont
		\textbf{What is irreducible error in the bias-variance decomposition?}
		
		\begin{enumerate}
			\item Error that can be reduced by increasing model complexity.
			\item Error that can be reduced by using more training data.
			\item Error that is due to noise in the data and cannot be reduced by any model.
			\item Error that is due to the choice of loss function.
		\end{enumerate}
		
		\vspace{0.05cm}
		\textcolor{myblue}{\textbf{Answer: C}}
		\vspace{0.05cm}
		\textcolor{myblue}{\textbf{Explanation:}} Irreducible error = noise in data (measurement errors, inherent randomness). Cannot be reduced by any model regardless of complexity.
		
		\vspace{0.05cm}
		\textcolor{myred}{\textbf{Why Others Are Wrong:}}
		\begin{itemize}
			\item \textcolor{myred}{A:} This describes reducible error (bias) - can be reduced.
			\item \textcolor{myred}{B:} This describes reducible error (variance) - can be reduced.
			\item \textcolor{myred}{D:} Loss function choice is a modeling decision, not irreducible error.
		\end{itemize}
		
		\textcolor{myred}{\textbf{Exam Trap:}} Irreducible error = \textcolor{myblue}{noise in the data} - cannot be eliminated! Even the best possible model has irreducible error.
	\end{frame}
	
	\begin{frame}{Q60: Bias-Variance Tradeoff - Summary}
		\fontsize{6.5}{7.5}\selectfont
		\textbf{Which of the following statements about the bias-variance tradeoff is TRUE?}
		
		\begin{enumerate}
			\item Simple models tend to have low bias and high variance.
			\item Complex models tend to have high bias and low variance.
			\item The optimal model balances bias and variance to minimize total error.
			\item Bias and variance are independent of each other.
		\end{enumerate}
		
		\vspace{0.05cm}
		\textcolor{myblue}{\textbf{Answer: C}}
		\vspace{0.05cm}
		\textcolor{myblue}{\textbf{Explanation:}} Optimal model balances bias and variance to minimize total error (bias$^2$ + variance + irreducible error).
		
		\vspace{0.05cm}
		\textcolor{myred}{\textbf{Why Others Are Wrong:}}
		\begin{itemize}
			\item \textcolor{myred}{A:} False - simple models have high bias, low variance.
			\item \textcolor{myred}{B:} False - complex models have low bias, high variance.
			\item \textcolor{myred}{D:} False - bias and variance are related through model complexity.
		\end{itemize}
		
		\textcolor{myred}{\textbf{Exam Trap:}} Optimal model \textcolor{myblue}{balances bias and variance}! Not at extremes, and not independent.
	\end{frame}
	
	% ============================================================
	% SECTION 5: REGULARIZATION - QUESTIONS 61-75
	% ============================================================
	
	\begin{frame}{Q61: Regularization - Definition}
		\fontsize{6.5}{7.5}\selectfont
		\textbf{What is regularization in supervised learning?}
		
		\begin{enumerate}
			\item A technique to increase model complexity.
			\item A technique to reduce overfitting by adding a penalty term to the loss function.
			\item A technique to increase the training data size.
			\item A technique to select the best features.
		\end{enumerate}
		
		\vspace{0.05cm}
		\textcolor{myblue}{\textbf{Answer: B}}
		\vspace{0.05cm}
		\textcolor{myblue}{\textbf{Explanation:}} Regularization prevents overfitting by adding penalty term to loss function that discourages large parameter values or complex models.
		
		\vspace{0.05cm}
		\textcolor{myred}{\textbf{Why Others Are Wrong:}}
		\begin{itemize}
			\item \textcolor{myred}{A:} This is the opposite of regularization (which reduces complexity).
			\item \textcolor{myred}{C:} This describes data augmentation - a different technique.
			\item \textcolor{myred}{D:} This describes feature selection - related but distinct.
		\end{itemize}
		
		\textcolor{myred}{\textbf{Exam Trap:}} Regularization adds a \textcolor{myblue}{penalty} to the loss function! Not increasing complexity or data.
	\end{frame}
	
	\begin{frame}{Q62: L1 Regularization (Lasso)}
		\fontsize{6.5}{7.5}\selectfont
		\textbf{What is L1 regularization (Lasso) and what is its key property?}
		
		\begin{enumerate}
			\item Adds $\lambda \sum w_i^2$ to the loss; shrinks weights toward zero.
			\item Adds $\lambda \sum |w_i|$ to the loss; can set weights to exactly zero (feature selection).
			\item Adds $\lambda \sum w_i$ to the loss; increases weights.
			\item Adds $\lambda \sum |w_i|^3$ to the loss; has no special property.
		\end{enumerate}
		
		\vspace{0.05cm}
		\textcolor{myblue}{\textbf{Answer: B}}
		\vspace{0.05cm}
		\textcolor{myblue}{\textbf{Explanation:}} L1 adds $\lambda \sum |w_i|$ to loss. Key property: can set weights to exactly zero, performing feature selection. Also known as Lasso.
		
		\vspace{0.05cm}
		\textcolor{myred}{\textbf{Why Others Are Wrong:}}
		\begin{itemize}
			\item \textcolor{myred}{A:} This describes L2 regularization (Ridge) - shrinks but doesn't zero out.
			\item \textcolor{myred}{C:} Linear penalty is not standard and would increase weights.
			\item \textcolor{myred}{D:} Cubic penalty is not a standard regularization technique.
		\end{itemize}
		
		\textcolor{myred}{\textbf{Exam Trap:}} L1 regularization \textcolor{myblue}{can set weights to zero} (feature selection)! This is the key differentiator from L2.
	\end{frame}
	
	\begin{frame}{Q63: L2 Regularization (Ridge)}
		\fontsize{6.5}{7.5}\selectfont
		\textbf{What is L2 regularization (Ridge) and what is its key property?}
		
		\begin{enumerate}
			\item Adds $\lambda \sum |w_i|$ to the loss; can set weights to zero.
			\item Adds $\lambda \sum w_i^2$ to the loss; shrinks weights toward zero but does not set them to zero.
			\item Adds $\lambda \sum w_i$ to the loss; increases weights.
			\item Adds $\lambda \sum |w_i|^3$ to the loss; has no special property.
		\end{enumerate}
		
		\vspace{0.05cm}
		\textcolor{myblue}{\textbf{Answer: B}}
		\vspace{0.05cm}
		\textcolor{myblue}{\textbf{Explanation:}} L2 adds $\lambda \sum w_i^2$ to loss. Shrinks weights toward zero but doesn't set them to zero. Also known as Ridge regression.
		
		\vspace{0.05cm}
		\textcolor{myred}{\textbf{Why Others Are Wrong:}}
		\begin{itemize}
			\item \textcolor{myred}{A:} This describes L1 regularization (Lasso) - can set to zero.
			\item \textcolor{myred}{C:} Linear penalty is not standard regularization.
			\item \textcolor{myred}{D:} Cubic penalty is not standard regularization.
		\end{itemize}
		
		\textcolor{myred}{\textbf{Exam Trap:}} L2 regularization \textcolor{myblue}{shrinks weights} but does not set them to zero! This is the key differentiator from L1.
	\end{frame}
	
	\begin{frame}{Q64: Elastic Net}
		\fontsize{6.5}{7.5}\selectfont
		\textbf{What is Elastic Net regularization?}
		
		\begin{enumerate}
			\item A combination of L1 and L2 regularization: $\lambda_1 \sum |w_i| + \lambda_2 \sum w_i^2$
			\item A method that only uses L1 regularization.
			\item A method that only uses L2 regularization.
			\item A method that adds no regularization.
		\end{enumerate}
		
		\vspace{0.05cm}
		\textcolor{myblue}{\textbf{Answer: A}}
		\vspace{0.05cm}
		\textcolor{myblue}{\textbf{Explanation:}} Elastic Net combines L1 and L2: $\lambda_1 \sum |w_i| + \lambda_2 \sum w_i^2$. Benefits from both feature selection (L1) and weight shrinkage (L2).
		
		\vspace{0.05cm}
		\textcolor{myred}{\textbf{Why Others Are Wrong:}}
		\begin{itemize}
			\item \textcolor{myred}{B:} This is Lasso (L1 only) - incomplete.
			\item \textcolor{myred}{C:} This is Ridge (L2 only) - incomplete.
			\item \textcolor{myred}{D:} This is no regularization - not Elastic Net.
		\end{itemize}
		
		\textcolor{myred}{\textbf{Exam Trap:}} Elastic Net = \textcolor{myblue}{L1 + L2} regularization! It combines both penalties.
	\end{frame}
	
	\begin{frame}{Q65: Regularization Parameter $\lambda$}
		\fontsize{6.5}{7.5}\selectfont
		\textbf{What happens as the regularization parameter $\lambda$ increases?}
		
		\begin{enumerate}
			\item The model becomes more complex.
			\item The model becomes simpler (weights shrink more toward zero).
			\item The model becomes more prone to overfitting.
			\item The model has no change.
		\end{enumerate}
		
		\vspace{0.05cm}
		\textcolor{myblue}{\textbf{Answer: B}}
		\vspace{0.05cm}
		\textcolor{myblue}{\textbf{Explanation:}} As $\lambda$ increases, penalty on large weights becomes stronger, weights shrink more toward zero. Model becomes simpler (helps prevent overfitting).
		
		\vspace{0.05cm}
		\textcolor{myred}{\textbf{Why Others Are Wrong:}}
		\begin{itemize}
			\item \textcolor{myred}{A:} Larger $\lambda$ means more regularization = simpler model, not more complex.
			\item \textcolor{myred}{C:} Larger $\lambda$ reduces overfitting, doesn't increase it.
			\item \textcolor{myred}{D:} $\lambda$ definitely changes the model - larger = simpler.
		\end{itemize}
		
		\textcolor{myred}{\textbf{Exam Trap:}} Larger $\lambda$ $\rightarrow$ \textcolor{myblue}{simpler model} (more regularization)! Not more complex.
	\end{frame}
	
	\begin{frame}{Q66: Regularization - Overfitting Prevention}
		\fontsize{6.5}{7.5}\selectfont
		\textbf{How does regularization prevent overfitting?}
		
		\begin{enumerate}
			\item By increasing the number of features.
			\item By penalizing large parameter values, which reduces the model's complexity and sensitivity to training data.
			\item By increasing the training data size.
			\item By using a more complex loss function.
		\end{enumerate}
		
		\vspace{0.05cm}
		\textcolor{myblue}{\textbf{Answer: B}}
		\vspace{0.05cm}
		\textcolor{myblue}{\textbf{Explanation:}} Regularization penalizes large parameter values. Reduces model complexity and sensitivity to training data fluctuations, improving generalization.
		
		\vspace{0.05cm}
		\textcolor{myred}{\textbf{Why Others Are Wrong:}}
		\begin{itemize}
			\item \textcolor{myred}{A:} Increasing features increases overfitting risk, doesn't prevent it.
			\item \textcolor{myred}{C:} This describes data augmentation, not regularization.
			\item \textcolor{myred}{D:} More complex loss function isn't regularization.
		\end{itemize}
		
		\textcolor{myred}{\textbf{Exam Trap:}} Regularization = \textcolor{myblue}{penalizes large weights} to reduce overfitting! This is the mechanism.
	\end{frame}
	
	\begin{frame}{Q67: Dropout as Regularization}
		\fontsize{6.5}{7.5}\selectfont
		\textbf{Is dropout a form of regularization? If so, how does it work?}
		
		\begin{enumerate}
			\item No, dropout is not a form of regularization.
			\item Yes, dropout randomly drops neurons during training, creating an ensemble effect and preventing co-adaptation.
			\item Yes, dropout adds a penalty term to the loss function.
			\item No, dropout only affects the forward pass, not the backward pass.
		\end{enumerate}
		
		\vspace{0.05cm}
		\textcolor{myblue}{\textbf{Answer: B}}
		\vspace{0.05cm}
		\textcolor{myblue}{\textbf{Explanation:}} Dropout is a regularization technique (and ensemble technique). Randomly drops (sets to zero) a fraction of neurons during each training iteration. Prevents co-adaptation.
		
		\vspace{0.05cm}
		\textcolor{myred}{\textbf{Why Others Are Wrong:}}
		\begin{itemize}
			\item \textcolor{myred}{A:} False - dropout IS a form of regularization.
			\item \textcolor{myred}{C:} Dropout doesn't add a penalty term - it's a different mechanism.
			\item \textcolor{myred}{D:} Dropout affects both forward and backward passes.
		\end{itemize}
		
		\textcolor{myred}{\textbf{Exam Trap:}} Dropout is a \textcolor{myblue}{regularization technique} for neural networks! It works by randomly dropping neurons, not by adding penalty.
	\end{frame}
	
	\begin{frame}{Q68: Early Stopping as Regularization}
		\fontsize{6.5}{7.5}\selectfont
		\textbf{Is early stopping a form of regularization? If so, how does it work?}
		
		\begin{enumerate}
			\item No, early stopping is not a form of regularization.
			\item Yes, early stopping stops training before convergence, limiting the model's complexity and preventing overfitting.
			\item Yes, early stopping adds a penalty term to the loss function.
			\item No, early stopping only affects the training speed, not the model.
		\end{enumerate}
		
		\vspace{0.05cm}
		\textcolor{myblue}{\textbf{Answer: B}}
		\vspace{0.05cm}
		\textcolor{myblue}{\textbf{Explanation:}} Early stopping is a regularization technique. Stops training when validation error starts to increase, limiting model complexity and preventing overfitting.
		
		\vspace{0.05cm}
		\textcolor{myred}{\textbf{Why Others Are Wrong:}}
		\begin{itemize}
			\item \textcolor{myred}{A:} False - early stopping IS a form of regularization.
			\item \textcolor{myred}{C:} Early stopping doesn't add penalty term - different mechanism.
			\item \textcolor{myred}{D:} Early stopping affects model (stops before overfitting).
		\end{itemize}
		
		\textcolor{myred}{\textbf{Exam Trap:}} Early stopping is a \textcolor{myblue}{regularization technique} that stops training early! It's a form of "implicit regularization."
	\end{frame}
	
	\begin{frame}{Q69: Regularization - Bias-Variance Tradeoff}
		\fontsize{6.5}{7.5}\selectfont
		\textbf{How does regularization affect the bias-variance tradeoff?}
		
		\begin{enumerate}
			\item Regularization increases both bias and variance.
			\item Regularization decreases bias and increases variance.
			\item Regularization increases bias and decreases variance.
			\item Regularization has no effect on bias or variance.
		\end{enumerate}
		
		\vspace{0.05cm}
		\textcolor{myblue}{\textbf{Answer: C}}
		\vspace{0.05cm}
		\textcolor{myblue}{\textbf{Explanation:}} Regularization increases bias (simpler model may not fit as well) and decreases variance (less sensitive to training data). This is the tradeoff.
		
		\vspace{0.05cm}
		\textcolor{myred}{\textbf{Why Others Are Wrong:}}
		\begin{itemize}
			\item \textcolor{myred}{A:} Both increase is incorrect - variance decreases.
			\item \textcolor{myred}{B:} This is the opposite direction (regularization $\to$ bias $\uparrow$, variance $\downarrow$).
			\item \textcolor{myred}{D:} False - regularization definitely affects both.
		\end{itemize}
		
		\textcolor{myred}{\textbf{Exam Trap:}} Regularization $\rightarrow$ \textcolor{myblue}{increases bias, decreases variance}! This is how it prevents overfitting.
	\end{frame}
	
	\begin{frame}{Q70: Regularization - When to Use}
		\fontsize{6.5}{7.5}\selectfont
		\textbf{When should regularization be applied in supervised learning?}
		
		\begin{enumerate}
			\item When the model is underfitting.
			\item When the model is overfitting (high variance).
			\item When the model has too few parameters.
			\item When the training data is perfectly clean.
		\end{enumerate}
		
		\vspace{0.05cm}
		\textcolor{myblue}{\textbf{Answer: B}}
		\vspace{0.05cm}
		\textcolor{myblue}{\textbf{Explanation:}} Regularization should be applied when model is overfitting (high variance). It reduces overfitting by penalizing large weights and limiting complexity.
		
		\vspace{0.05cm}
		\textcolor{myred}{\textbf{Why Others Are Wrong:}}
		\begin{itemize}
			\item \textcolor{myred}{A:} Underfitting would be worsened by regularization (makes model simpler).
			\item \textcolor{myred}{C:} Too few parameters = underfitting - regularization would make worse.
			\item \textcolor{myred}{D:} Perfectly clean data doesn't need regularization (though still might help).
		\end{itemize}
		
		\textcolor{myred}{\textbf{Exam Trap:}} Use regularization when \textcolor{myblue}{overfitting} (high variance)! Not when underfitting.
	\end{frame}
	
	\begin{frame}{Q71: L1 vs L2 - Feature Selection}
		\fontsize{6.5}{7.5}\selectfont
		\textbf{Which regularization method is more likely to perform feature selection?}
		
		\begin{enumerate}
			\item L2 regularization (Ridge)
			\item L1 regularization (Lasso)
			\item Both perform feature selection equally.
			\item Neither performs feature selection.
		\end{enumerate}
		
		\vspace{0.05cm}
		\textcolor{myblue}{\textbf{Answer: B}}
		\vspace{0.05cm}
		\textcolor{myblue}{\textbf{Explanation:}} L1 (Lasso) performs feature selection because it can set weights to exactly zero. L2 (Ridge) shrinks weights but doesn't zero them out.
		
		\vspace{0.05cm}
		\textcolor{myred}{\textbf{Why Others Are Wrong:}}
		\begin{itemize}
			\item \textcolor{myred}{A:} L2 only shrinks weights, doesn't zero them - no feature selection.
			\item \textcolor{myred}{C:} They are not equal - L1 does feature selection, L2 doesn't.
			\item \textcolor{myred}{D:} L1 definitely performs feature selection.
		\end{itemize}
		
		\textcolor{myred}{\textbf{Exam Trap:}} L1 (Lasso) = \textcolor{myblue}{feature selection}; L2 (Ridge) = \textcolor{myblue}{weight shrinkage}! This is a key distinction.
	\end{frame}
	
	\begin{frame}{Q72: Regularization - Interpretation}
		\fontsize{6.5}{7.5}\selectfont
		\textbf{In the context of regularization, what does a weight of zero indicate?}
		
		\begin{enumerate}
			\item The feature is very important.
			\item The feature is not used by the model (feature selection).
			\item The feature has high variance.
			\item The feature has high bias.
		\end{enumerate}
		
		\vspace{0.05cm}
		\textcolor{myblue}{\textbf{Answer: B}}
		\vspace{0.05cm}
		\textcolor{myblue}{\textbf{Explanation:}} Weight of zero indicates feature is not used by model. This is feature selection (especially from L1 regularization).
		
		\vspace{0.05cm}
		\textcolor{myred}{\textbf{Why Others Are Wrong:}}
		\begin{itemize}
			\item \textcolor{myred}{A:} Weight of zero means feature has no importance.
			\item \textcolor{myred}{C:} Weight value doesn't directly indicate variance.
			\item \textcolor{myred}{D:} Weight value doesn't directly indicate bias.
		\end{itemize}
		
		\textcolor{myred}{\textbf{Exam Trap:}} Weight = 0 $\rightarrow$ \textcolor{myblue}{feature not used} (feature selection)! This is the interpretation.
	\end{frame}
	
	\begin{frame}{Q73: Regularization - Overfitting vs Underfitting}
		\fontsize{6.5}{7.5}\selectfont
		\textbf{If $\lambda$ is too large in regularization, what happens?}
		
		\begin{enumerate}
			\item The model overfits.
			\item The model underfits.
			\item The model performs optimally.
			\item The model has no effect.
		\end{enumerate}
		
		\vspace{0.05cm}
		\textcolor{myblue}{\textbf{Answer: B}}
		\vspace{0.05cm}
		\textcolor{myblue}{\textbf{Explanation:}} If $\lambda$ too large, penalty too strong, all weights shrink to zero. Model becomes too simple and underfits.
		
		\vspace{0.05cm}
		\textcolor{myred}{\textbf{Why Others Are Wrong:}}
		\begin{itemize}
			\item \textcolor{myred}{A:} Too small $\lambda$ leads to overfitting, not too large.
			\item \textcolor{myred}{C:} Optimal requires balanced $\lambda$ - not too large or small.
			\item \textcolor{myred}{D:} $\lambda$ definitely affects model - too large causes underfitting.
		\end{itemize}
		
		\textcolor{myred}{\textbf{Exam Trap:}} Too large $\lambda$ $\rightarrow$ \textcolor{myred}{underfitting}! Too small $\lambda$ $\rightarrow$ overfitting. Need balance.
	\end{frame}
	
	\begin{frame}{Q74: Regularization - Summary}
		\fontsize{6.5}{7.5}\selectfont
		\textbf{Which of the following is a form of regularization?}
		
		\begin{enumerate}
			\item L1 regularization (Lasso)
			\item L2 regularization (Ridge)
			\item Dropout
			\item All of the above
		\end{enumerate}
		
		\vspace{0.05cm}
		\textcolor{myblue}{\textbf{Answer: D}}
		\vspace{0.05cm}
		\textcolor{myblue}{\textbf{Explanation:}} L1, L2, and dropout are all forms of regularization. They prevent overfitting through different mechanisms.
		
		\vspace{0.05cm}
		\textcolor{myred}{\textbf{Why Others Are Wrong:}}
		\begin{itemize}
			\item \textcolor{myred}{A:} Correct but incomplete - other options also valid.
			\item \textcolor{myred}{B:} Correct but incomplete - other options also valid.
			\item \textcolor{myred}{C:} Correct but incomplete - other options also valid.
		\end{itemize}
		
		\textcolor{myred}{\textbf{Exam Trap:}} Many techniques are forms of regularization! Don't assume only L1/L2 - dropout and early stopping also count.
	\end{frame}
	
	\begin{frame}{Q75: Regularization - Summary}
		\fontsize{6.5}{7.5}\selectfont
		\textbf{What is the purpose of the regularization parameter $\lambda$?}
		
		\begin{enumerate}
			\item To control the tradeoff between fitting the training data and keeping the model simple.
			\item To control the learning rate.
			\item To control the number of epochs.
			\item To control the batch size.
		\end{enumerate}
		
		\vspace{0.05cm}
		\textcolor{myblue}{\textbf{Answer: A}}
		\vspace{0.05cm}
		\textcolor{myblue}{\textbf{Explanation:}} $\lambda$ controls tradeoff between fitting training data (small $\lambda$) and keeping model simple (large $\lambda$).
		
		\vspace{0.05cm}
		\textcolor{myred}{\textbf{Why Others Are Wrong:}}
		\begin{itemize}
			\item \textcolor{myred}{B:} Learning rate is $\alpha$ - different hyperparameter.
			\item \textcolor{myred}{C:} Number of epochs is a different hyperparameter.
			\item \textcolor{myred}{D:} Batch size is a different hyperparameter.
		\end{itemize}
		
		\textcolor{myred}{\textbf{Exam Trap:}} $\lambda$ controls the \textcolor{myblue}{bias-variance tradeoff} in regularization! Not learning rate, epochs, or batch size.
	\end{frame}
	
	% ============================================================
	% SECTION 6: CROSS-VALIDATION - QUESTIONS 76-90
	% ============================================================
	
	\begin{frame}{Q76: Cross-Validation - Definition}
		\fontsize{6.5}{7.5}\selectfont
		\textbf{What is cross-validation in supervised learning?}
		
		\begin{enumerate}
			\item A method to increase the training data size.
			\item A resampling technique used to evaluate model performance and select hyperparameters.
			\item A method to reduce overfitting by adding a penalty term.
			\item A method to initialize model parameters.
		\end{enumerate}
		
		\vspace{0.05cm}
		\textcolor{myblue}{\textbf{Answer: B}}
		\vspace{0.05cm}
		\textcolor{myblue}{\textbf{Explanation:}} Cross-validation partitions data into training and validation sets multiple times to get reliable performance estimates and select hyperparameters.
		
		\vspace{0.05cm}
		\textcolor{myred}{\textbf{Why Others Are Wrong:}}
		\begin{itemize}
			\item \textcolor{myred}{A:} This describes data augmentation - different technique.
			\item \textcolor{myred}{C:} This describes regularization - different technique.
			\item \textcolor{myred}{D:} This describes weight initialization - different concept.
		\end{itemize}
		
		\textcolor{myred}{\textbf{Exam Trap:}} Cross-validation evaluates \textcolor{myblue}{model performance} and selects \textcolor{myblue}{hyperparameters}! Not data augmentation or regularization.
	\end{frame}
	
	\begin{frame}{Q77: K-Fold Cross-Validation}
		\fontsize{6.5}{7.5}\selectfont
		\textbf{How does k-fold cross-validation work?}
		
		\begin{enumerate}
			\item Split data into k parts, train on one part and validate on the remaining k-1 parts.
			\item Split data into k parts, train on k-1 parts and validate on the remaining 1 part, repeat k times.
			\item Split data into k parts and use all parts for training.
			\item Split data into k parts and use all parts for validation.
		\end{enumerate}
		
		\vspace{0.05cm}
		\textcolor{myblue}{\textbf{Answer: B}}
		\vspace{0.05cm}
		\textcolor{myblue}{\textbf{Explanation:}} Split into k folds. For each iteration: train on k-1 folds, validate on 1 fold. Repeat k times. Average results.
		
		\vspace{0.05cm}
		\textcolor{myred}{\textbf{Why Others Are Wrong:}}
		\begin{itemize}
			\item \textcolor{myred}{A:} Reverses training and validation (would train on too little data).
			\item \textcolor{myred}{C:} Using all data for training = no validation - not cross-validation.
			\item \textcolor{myred}{D:} Using all data for validation = no training - not cross-validation.
		\end{itemize}
		
		\textcolor{myred}{\textbf{Exam Trap:}} K-fold CV = train on \textcolor{myblue}{k-1 folds}, validate on \textcolor{myblue}{1 fold}, repeat k times! This is the exact procedure.
	\end{frame}
	
	\begin{frame}{Q78: Stratified Cross-Validation}
		\fontsize{6.5}{7.5}\selectfont
		\textbf{What is stratified cross-validation?}
		
		\begin{enumerate}
			\item Cross-validation where each fold has the same number of observations.
			\item Cross-validation where each fold preserves the same class distribution as the full dataset.
			\item Cross-validation where the folds are chosen randomly.
			\item Cross-validation where the folds are chosen sequentially.
		\end{enumerate}
		
		\vspace{0.05cm}
		\textcolor{myblue}{\textbf{Answer: B}}
		\vspace{0.05cm}
		\textcolor{myblue}{\textbf{Explanation:}} Stratified CV ensures each fold has same class distribution (proportion of each class) as the full dataset. Important for imbalanced datasets.
		
		\vspace{0.05cm}
		\textcolor{myred}{\textbf{Why Others Are Wrong:}}
		\begin{itemize}
			\item \textcolor{myred}{A:} Equal fold sizes is standard k-fold, not stratification.
			\item \textcolor{myred}{C:} Random folds may not preserve class distribution.
			\item \textcolor{myred}{D:} Sequential folds may introduce order bias.
		\end{itemize}
		
		\textcolor{myred}{\textbf{Exam Trap:}} Stratified CV preserves \textcolor{myblue}{class proportions} in each fold! This is the key differentiator.
	\end{frame}
	
	\begin{frame}{Q79: Stratified CV - When to Use}
		\fontsize{6.5}{7.5}\selectfont
		\textbf{When is stratified cross-validation particularly important?}
		
		\begin{enumerate}
			\item When the dataset is perfectly balanced.
			\item When the dataset is imbalanced (one class is much rarer than others).
			\item When the dataset is very large.
			\item When the dataset has no categorical variables.
		\end{enumerate}
		
		\vspace{0.05cm}
		\textcolor{myblue}{\textbf{Answer: B}}
		\vspace{0.05cm}
		\textcolor{myblue}{\textbf{Explanation:}} Stratified CV is especially important for imbalanced datasets. Without it, some folds might have very few samples from the rare class.
		
		\vspace{0.05cm}
		\textcolor{myred}{\textbf{Why Others Are Wrong:}}
		\begin{itemize}
			\item \textcolor{myred}{A:} Balanced datasets don't need stratification (though it doesn't hurt).
			\item \textcolor{myred}{C:} Dataset size doesn't determine need for stratification.
			\item \textcolor{myred}{D:} Categorical variables not relevant - stratification is about class distribution.
		\end{itemize}
		
		\textcolor{myred}{\textbf{Exam Trap:}} Stratified CV is essential for \textcolor{myblue}{imbalanced datasets}! This is the primary use case.
	\end{frame}
	
	\begin{frame}{Q80: Leave-One-Out Cross-Validation (LOOCV)}
		\fontsize{6.5}{7.5}\selectfont
		\textbf{What is Leave-One-Out Cross-Validation (LOOCV)?}
		
		\begin{enumerate}
			\item A special case of k-fold CV where k = N (one observation per fold).
			\item A special case of k-fold CV where k = 2.
			\item A special case of k-fold CV where k = 10.
			\item A method that uses no validation set.
		\end{enumerate}
		
		\vspace{0.05cm}
		\textcolor{myblue}{\textbf{Answer: A}}
		\vspace{0.05cm}
		\textcolor{myblue}{\textbf{Explanation:}} LOOCV = k-fold CV with k = N (number of observations). Each fold has one observation. Train on N-1, validate on 1. Repeat N times.
		
		\vspace{0.05cm}
		\textcolor{myred}{\textbf{Why Others Are Wrong:}}
		\begin{itemize}
			\item \textcolor{myred}{B:} k=2 is 2-fold CV - not LOOCV (unless N=2).
			\item \textcolor{myred}{C:} k=10 is 10-fold CV - most common but not LOOCV.
			\item \textcolor{myred}{D:} LOOCV uses validation sets (one observation at a time).
		\end{itemize}
		
		\textcolor{myred}{\textbf{Exam Trap:}} LOOCV = \textcolor{myblue}{k-fold CV with k = N}! This is the precise definition.
	\end{frame}
	
	\begin{frame}{Q81: LOOCV - Advantages}
		\fontsize{6.5}{7.5}\selectfont
		\textbf{What is an advantage of Leave-One-Out Cross-Validation?}
		
		\begin{enumerate}
			\item It is computationally cheap.
			\item It uses almost all data for training each time (N-1 observations), giving low bias.
			\item It always has low variance.
			\item It works well for very large datasets.
		\end{enumerate}
		
		\vspace{0.05cm}
		\textcolor{myblue}{\textbf{Answer: B}}
		\vspace{0.05cm}
		\textcolor{myblue}{\textbf{Explanation:}} LOOCV uses N-1 observations for training each time. This gives low bias (almost full dataset used for training).
		
		\vspace{0.05cm}
		\textcolor{myred}{\textbf{Why Others Are Wrong:}}
		\begin{itemize}
			\item \textcolor{myred}{A:} LOOCV is computationally expensive (requires N training runs).
			\item \textcolor{myred}{C:} LOOCV can have high variance (estimates are correlated).
			\item \textcolor{myred}{D:} LOOCV is not suitable for large datasets (too expensive).
		\end{itemize}
		
		\textcolor{myred}{\textbf{Exam Trap:}} LOOCV has \textcolor{myblue}{low bias} but is \textcolor{myred}{computationally expensive}! This is the key tradeoff.
	\end{frame}
	
	\begin{frame}{Q82: LOOCV - Disadvantages}
		\fontsize{6.5}{7.5}\selectfont
		\textbf{What is a disadvantage of Leave-One-Out Cross-Validation?}
		
		\begin{enumerate}
			\item It has high bias.
			\item It is computationally expensive (requires N training runs).
			\item It always overfits.
			\item It cannot be used for classification.
		\end{enumerate}
		
		\vspace{0.05cm}
		\textcolor{myblue}{\textbf{Answer: B}}
		\vspace{0.05cm}
		\textcolor{myblue}{\textbf{Explanation:}} LOOCV requires training model N times (once for each held-out observation). This is impractical for large datasets.
		
		\vspace{0.05cm}
		\textcolor{myred}{\textbf{Why Others Are Wrong:}}
		\begin{itemize}
			\item \textcolor{myred}{A:} LOOCV has low bias (uses N-1 training samples).
			\item \textcolor{myred}{C:} LOOCV doesn't necessarily overfit (though variance can be high).
			\item \textcolor{myred}{D:} LOOCV can be used for classification too.
		\end{itemize}
		
		\textcolor{myred}{\textbf{Exam Trap:}} LOOCV is \textcolor{myred}{computationally expensive} for large N! This is the main disadvantage.
	\end{frame}
	
	\begin{frame}{Q83: Cross-Validation - Bias-Variance Tradeoff}
		\fontsize{6.5}{7.5}\selectfont
		\textbf{How does the choice of k in k-fold cross-validation affect the bias-variance tradeoff?}
		
		\begin{enumerate}
			\item Larger k increases bias and decreases variance.
			\item Larger k decreases bias and increases variance.
			\item Larger k has no effect on bias or variance.
			\item Smaller k decreases bias and increases variance.
		\end{enumerate}
		
		\vspace{0.05cm}
		\textcolor{myblue}{\textbf{Answer: B}}
		\vspace{0.05cm}
		\textcolor{myblue}{\textbf{Explanation:}} Larger k = more training data (k-1 folds) $\to$ lower bias, higher variance. LOOCV (k=N) has lowest bias, highest variance.
		
		\vspace{0.05cm}
		\textcolor{myred}{\textbf{Why Others Are Wrong:}}
		\begin{itemize}
			\item \textcolor{myred}{A:} This is opposite - larger k decreases bias.
			\item \textcolor{myred}{C:} False - k definitely affects bias-variance.
			\item \textcolor{myred}{D:} Smaller k increases bias, decreases variance.
		\end{itemize}
		
		\textcolor{myred}{\textbf{Exam Trap:}} Larger k $\rightarrow$ \textcolor{myblue}{lower bias, higher variance}! This is the tradeoff.
	\end{frame}
	
	\begin{frame}{Q84: Cross-Validation - Common k Value}
		\fontsize{6.5}{7.5}\selectfont
		\textbf{What is the most common value of k used in k-fold cross-validation?}
		
		\begin{enumerate}
			\item k = 2
			\item k = 5
			\item k = 10
			\item k = N (LOOCV)
		\end{enumerate}
		
		\vspace{0.05cm}
		\textcolor{myblue}{\textbf{Answer: C}}
		\vspace{0.05cm}
		\textcolor{myblue}{\textbf{Explanation:}} k = 10 is most common. Provides good balance between bias and variance and is computationally feasible.
		
		\vspace{0.05cm}
		\textcolor{myred}{\textbf{Why Others Are Wrong:}}
		\begin{itemize}
			\item \textcolor{myred}{A:} k=2 is sometimes used but not most common.
			\item \textcolor{myred}{B:} k=5 is also common but k=10 is most common.
			\item \textcolor{myred}{D:} LOOCV is used for small datasets, not most common.
		\end{itemize}
		
		\textcolor{myred}{\textbf{Exam Trap:}} k = 10 is the \textcolor{myblue}{most common} value! This is a standard convention.
	\end{frame}
	
	\begin{frame}{Q85: Cross-Validation - Purpose}
		\fontsize{6.5}{7.5}\selectfont
		\textbf{What is the primary purpose of cross-validation?}
		
		\begin{enumerate}
			\item To increase the training data size.
			\item To evaluate how well a model will generalize to new, unseen data.
			\item To reduce the number of features.
			\item To increase the learning rate.
		\end{enumerate}
		
		\vspace{0.05cm}
		\textcolor{myblue}{\textbf{Answer: B}}
		\vspace{0.05cm}
		\textcolor{myblue}{\textbf{Explanation:}} Primary purpose: evaluate model generalization performance. Provides reliable estimate of test performance.
		
		\vspace{0.05cm}
		\textcolor{myred}{\textbf{Why Others Are Wrong:}}
		\begin{itemize}
			\item \textcolor{myred}{A:} Data augmentation increases data size - not CV.
			\item \textcolor{myred}{C:} Feature selection reduces features - not CV.
			\item \textcolor{myred}{D:} Learning rate is a training hyperparameter - not CV.
		\end{itemize}
		
		\textcolor{myred}{\textbf{Exam Trap:}} Cross-validation estimates \textcolor{myblue}{generalization performance}! This is the fundamental purpose.
	\end{frame}
	
	\begin{frame}{Q86: Bootstrapping - Definition}
		\fontsize{6.5}{7.5}\selectfont
		\textbf{What is bootstrapping in the context of model evaluation?}
		
		\begin{enumerate}
			\item A resampling technique that samples with replacement from the dataset.
			\item A resampling technique that samples without replacement.
			\item A method to add regularization.
			\item A method to reduce the number of features.
		\end{enumerate}
		
		\vspace{0.05cm}
		\textcolor{myblue}{\textbf{Answer: A}}
		\vspace{0.05cm}
		\textcolor{myblue}{\textbf{Explanation:}} Bootstrapping samples with replacement from dataset. Creates multiple bootstrap samples (same size as original) with some observations repeated and some omitted.
		
		\vspace{0.05cm}
		\textcolor{myred}{\textbf{Why Others Are Wrong:}}
		\begin{itemize}
			\item \textcolor{myred}{B:} Sampling without replacement is subsampling - not bootstrapping.
			\item \textcolor{myred}{C:} Bootstrapping is not regularization.
			\item \textcolor{myred}{D:} Bootstrapping is not feature selection.
		\end{itemize}
		
		\textcolor{myred}{\textbf{Exam Trap:}} Bootstrapping = \textcolor{myblue}{sampling with replacement}! This is the key defining feature.
	\end{frame}
	
	\begin{frame}{Q87: Bootstrapping - Purpose}
		\fontsize{6.5}{7.5}\selectfont
		\textbf{What is the primary purpose of bootstrapping in model evaluation?}
		
		\begin{enumerate}
			\item To estimate the variability of model performance (e.g., confidence intervals).
			\item To reduce overfitting.
			\item To increase the learning rate.
			\item To select the best features.
		\end{enumerate}
		
		\vspace{0.05cm}
		\textcolor{myblue}{\textbf{Answer: A}}
		\vspace{0.05cm}
		\textcolor{myblue}{\textbf{Explanation:}} Bootstrapping estimates performance variability and confidence intervals. Trains multiple models on different bootstrap samples to get performance distribution.
		
		\vspace{0.05cm}
		\textcolor{myred}{\textbf{Why Others Are Wrong:}}
		\begin{itemize}
			\item \textcolor{myred}{B:} Bootstrapping is not regularization (though related to ensemble methods).
			\item \textcolor{myred}{C:} Learning rate is a training parameter - not bootstrapping.
			\item \textcolor{myred}{D:} Feature selection is a different technique.
		\end{itemize}
		
		\textcolor{myred}{\textbf{Exam Trap:}} Bootstrapping estimates \textcolor{myblue}{performance variability}! This is the primary use.
	\end{frame}
	
	\begin{frame}{Q88: Cross-Validation vs Bootstrapping}
		\fontsize{6.5}{7.5}\selectfont
		\textbf{What is the key difference between cross-validation and bootstrapping?}
		
		\begin{enumerate}
			\item Cross-validation samples without replacement; bootstrapping samples with replacement.
			\item Cross-validation uses all data for training; bootstrapping uses only part of the data.
			\item Cross-validation is for classification; bootstrapping is for regression.
			\item Cross-validation and bootstrapping are identical.
		\end{enumerate}
		
		\vspace{0.05cm}
		\textcolor{myblue}{\textbf{Answer: A}}
		\vspace{0.05cm}
		\textcolor{myblue}{\textbf{Explanation:}} CV (especially k-fold) samples without replacement (each observation appears in exactly one validation fold). Bootstrapping samples with replacement.
		
		\vspace{0.05cm}
		\textcolor{myred}{\textbf{Why Others Are Wrong:}}
		\begin{itemize}
			\item \textcolor{myred}{B:} False - both use parts of data for training/validation.
			\item \textcolor{myred}{C:} False - both can be used for either task.
			\item \textcolor{myred}{D:} False - they are different techniques.
		\end{itemize}
		
		\textcolor{myred}{\textbf{Exam Trap:}} CV = \textcolor{myblue}{without replacement}; Bootstrapping = \textcolor{myblue}{with replacement}! This is the key difference.
	\end{frame}
	
	\begin{frame}{Q89: Grid Search - Definition}
		\fontsize{6.5}{7.5}\selectfont
		\textbf{What is grid search in the context of hyperparameter tuning?}
		
		\begin{enumerate}
			\item A method that randomly searches for optimal hyperparameters.
			\item A method that exhaustively searches over a predefined grid of hyperparameter values.
			\item A method that uses gradient descent to find hyperparameters.
			\item A method that uses cross-validation to select the best model.
		\end{enumerate}
		
		\vspace{0.05cm}
		\textcolor{myblue}{\textbf{Answer: B}}
		\vspace{0.05cm}
		\textcolor{myblue}{\textbf{Explanation:}} Grid search exhaustively evaluates all combinations of hyperparameters from a predefined grid. Guarantees finding best within the grid but computationally expensive.
		
		\vspace{0.05cm}
		\textcolor{myred}{\textbf{Why Others Are Wrong:}}
		\begin{itemize}
			\item \textcolor{myred}{A:} This describes random search - different method.
			\item \textcolor{myred}{C:} This describes gradient-based hyperparameter optimization.
			\item \textcolor{myred}{D:} Cross-validation is used WITH grid search, not the same as grid search.
		\end{itemize}
		
		\textcolor{myred}{\textbf{Exam Trap:}} Grid search = \textcolor{myblue}{exhaustive search} over a parameter grid! This is the defining feature.
	\end{frame}
	
	\begin{frame}{Q90: Grid Search - Computational Cost}
		\fontsize{6.5}{7.5}\selectfont
		\textbf{What is a major disadvantage of grid search for hyperparameter tuning?}
		
		\begin{enumerate}
			\item It is computationally expensive for large grids.
			\item It always finds suboptimal hyperparameters.
			\item It cannot be used with cross-validation.
			\item It requires gradient information.
		\end{enumerate}
		
		\vspace{0.05cm}
		\textcolor{myblue}{\textbf{Answer: A}}
		\vspace{0.05cm}
		\textcolor{myblue}{\textbf{Explanation:}} Grid search is computationally expensive, especially for large grids. Number of combinations grows exponentially with number of hyperparameters.
		
		\vspace{0.05cm}
		\textcolor{myred}{\textbf{Why Others Are Wrong:}}
		\begin{itemize}
			\item \textcolor{myred}{B:} False - it finds best within the grid (not suboptimal).
			\item \textcolor{myred}{C:} False - grid search is often used with cross-validation.
			\item \textcolor{myred}{D:} False - grid search doesn't require gradients.
		\end{itemize}
		
		\textcolor{myred}{\textbf{Exam Trap:}} Grid search is \textcolor{myred}{computationally expensive} for large grids! This is the main limitation.
	\end{frame}
	
	% ============================================================
	% SECTION 7: ACCURACY-INTERPRETABILITY - QUESTIONS 91-100
	% ============================================================
	
	\begin{frame}{Q91: Accuracy-Interpretability Tradeoff}
		\fontsize{6.5}{7.5}\selectfont
		\textbf{What is the accuracy-interpretability tradeoff in supervised learning?}
		
		\begin{enumerate}
			\item The tradeoff between model complexity and training speed.
			\item The tradeoff between a model's predictive performance and how easily its decisions can be understood.
			\item The tradeoff between bias and variance.
			\item The tradeoff between the number of features and the number of observations.
		\end{enumerate}
		
		\vspace{0.05cm}
		\textcolor{myblue}{\textbf{Answer: B}}
		\vspace{0.05cm}
		\textcolor{myblue}{\textbf{Explanation:}} Accuracy-interpretability tradeoff: tension between predictive performance and understandability. More complex models often more accurate but less interpretable.
		
		\vspace{0.05cm}
		\textcolor{myred}{\textbf{Why Others Are Wrong:}}
		\begin{itemize}
			\item \textcolor{myred}{A:} This is speed-complexity tradeoff - different concept.
			\item \textcolor{myred}{C:} This is bias-variance tradeoff - different concept.
			\item \textcolor{myred}{D:} This is curse of dimensionality - different concept.
		\end{itemize}
		
		\textcolor{myred}{\textbf{Exam Trap:}} Accuracy vs Interpretability = \textcolor{myblue}{performance vs understandability}! This is the key concept.
	\end{frame}
	
	\begin{frame}{Q92: High Interpretability Models}
		\fontsize{6.5}{7.5}\selectfont
		\textbf{Which of the following models typically has high interpretability?}
		
		\begin{enumerate}
			\item Deep neural networks
			\item Random forests
			\item Decision trees (shallow)
			\item Gradient boosting machines
		\end{enumerate}
		
		\vspace{0.05cm}
		\textcolor{myblue}{\textbf{Answer: C}}
		\vspace{0.05cm}
		\textcolor{myblue}{\textbf{Explanation:}} Shallow decision trees have high interpretability. Decision-making process is easy to visualize and understand. "White-box" model.
		
		\vspace{0.05cm}
		\textcolor{myred}{\textbf{Why Others Are Wrong:}}
		\begin{itemize}
			\item \textcolor{myred}{A:} DNNs are black-box models (low interpretability).
			\item \textcolor{myred}{B:} Random forests are ensembles of trees (less interpretable).
			\item \textcolor{myred}{D:} GBM is complex ensemble (low interpretability).
		\end{itemize}
		
		\textcolor{myred}{\textbf{Exam Trap:}} \textcolor{myblue}{Decision trees} are highly interpretable "white-box" models! This is a key example.
	\end{frame}
	
	\begin{frame}{Q93: High Accuracy Models}
		\fontsize{6.5}{7.5}\selectfont
		\textbf{Which of the following models typically has high accuracy but low interpretability?}
		
		\begin{enumerate}
			\item Linear regression
			\item Decision trees (shallow)
			\item Deep neural networks
			\item Logistic regression
		\end{enumerate}
		
		\vspace{0.05cm}
		\textcolor{myblue}{\textbf{Answer: C}}
		\vspace{0.05cm}
		\textcolor{myblue}{\textbf{Explanation:}} Deep neural networks achieve high accuracy (especially on complex tasks) but have low interpretability. Often called "black-box" models.
		
		\vspace{0.05cm}
		\textcolor{myred}{\textbf{Why Others Are Wrong:}}
		\begin{itemize}
			\item \textcolor{myred}{A:} Linear regression is highly interpretable.
			\item \textcolor{myred}{B:} Shallow decision trees are highly interpretable.
			\item \textcolor{myred}{D:} Logistic regression is highly interpretable.
		\end{itemize}
		
		\textcolor{myred}{\textbf{Exam Trap:}} Deep neural networks = \textcolor{myblue}{high accuracy, low interpretability}! This is the classic example.
	\end{frame}
	
	\begin{frame}{Q94: Interpretability - Importance}
		\fontsize{6.5}{7.5}\selectfont
		\textbf{In which domain is interpretability particularly important?}
		
		\begin{enumerate}
			\item Image recognition
			\item Natural language processing
			\item Risk management and regulated industries (e.g., credit scoring)
			\item Recommendation systems
		\end{enumerate}
		
		\vspace{0.05cm}
		\textcolor{myblue}{\textbf{Answer: C}}
		\vspace{0.05cm}
		\textcolor{myblue}{\textbf{Explanation:}} Interpretability is critical in risk management and regulated industries (credit scoring, healthcare, insurance) where decisions must be explainable to regulators and customers.
		
		\vspace{0.05cm}
		\textcolor{myred}{\textbf{Why Others Are Wrong:}}
		\begin{itemize}
			\item \textcolor{myred}{A:} Image recognition prioritizes accuracy over interpretability.
			\item \textcolor{myred}{B:} NLP often prioritizes accuracy over interpretability.
			\item \textcolor{myred}{D:} Recommendation systems prioritize accuracy and personalization.
		\end{itemize}
		
		\textcolor{myred}{\textbf{Exam Trap:}} Interpretability is critical in \textcolor{myblue}{risk management and regulated industries}! This is the key domain.
	\end{frame}
	
	\begin{frame}{Q95: Risk Management - Model Choice}
		\fontsize{6.5}{7.5}\selectfont
		\textbf{In risk management, which type of model is often preferred despite potentially lower accuracy?}
		
		\begin{enumerate}
			\item Black-box models (e.g., deep neural networks)
			\item White-box models (e.g., decision trees, logistic regression)
			\item Ensemble models (e.g., random forests)
			\item Any model with the highest accuracy
		\end{enumerate}
		
		\vspace{0.05cm}
		\textcolor{myblue}{\textbf{Answer: B}}
		\vspace{0.05cm}
		\textcolor{myblue}{\textbf{Explanation:}} Risk management prefers white-box (interpretable) models like decision trees and logistic regression. Regulatory requirements demand transparency.
		
		\vspace{0.05cm}
		\textcolor{myred}{\textbf{Why Others Are Wrong:}}
		\begin{itemize}
			\item \textcolor{myred}{A:} Black-box models are avoided despite higher accuracy.
			\item \textcolor{myred}{C:} Ensembles are less interpretable than individual trees.
			\item \textcolor{myred}{D:} Highest accuracy ignores interpretability concerns.
		\end{itemize}
		
		\textcolor{myred}{\textbf{Exam Trap:}} Risk management prefers \textcolor{myblue}{interpretable (white-box)} models! This is the regulatory requirement.
	\end{frame}
	
	\begin{frame}{Q96: Accuracy-Interpretability Tradeoff - Summary}
		\fontsize{6.5}{7.5}\selectfont
		\textbf{Which of the following statements about the accuracy-interpretability tradeoff is TRUE?}
		
		\begin{enumerate}
			\item More interpretable models always have higher accuracy.
			\item Less interpretable models always have lower accuracy.
			\item There is often a tradeoff between accuracy and interpretability.
			\item Accuracy and interpretability are independent of each other.
		\end{enumerate}
		
		\vspace{0.05cm}
		\textcolor{myblue}{\textbf{Answer: C}}
		\vspace{0.05cm}
		\textcolor{myblue}{\textbf{Explanation:}} There is often a tradeoff. Simpler, more interpretable models may have lower accuracy; complex, less interpretable models may have higher accuracy.
		
		\vspace{0.05cm}
		\textcolor{myred}{\textbf{Why Others Are Wrong:}}
		\begin{itemize}
			\item \textcolor{myred}{A:} False - simpler models often have lower accuracy.
			\item \textcolor{myred}{B:} False - complex models often have higher accuracy.
			\item \textcolor{myred}{D:} False - they are related through model complexity.
		\end{itemize}
		
		\textcolor{myred}{\textbf{Exam Trap:}} Accuracy and interpretability are often \textcolor{myblue}{in tension} (tradeoff)! This is the key concept.
	\end{frame}
	
	\begin{frame}{Q97: Interpretability - Feature Importance}
		\fontsize{6.5}{7.5}\selectfont
		\textbf{What is feature importance in the context of model interpretability?}
		
		\begin{enumerate}
			\item A measure of how important each feature is in making predictions.
			\item A measure of how much data is available for each feature.
			\item A measure of the variance of each feature.
			\item A measure of the correlation between features.
		\end{enumerate}
		
		\vspace{0.05cm}
		\textcolor{myblue}{\textbf{Answer: A}}
		\vspace{0.05cm}
		\textcolor{myblue}{\textbf{Explanation:}} Feature importance measures how important each feature is in making predictions. Helps interpret models by identifying which features drive predictions.
		
		\vspace{0.05cm}
		\textcolor{myred}{\textbf{Why Others Are Wrong:}}
		\begin{itemize}
			\item \textcolor{myred}{B:} This is data availability - different concept.
			\item \textcolor{myred}{C:} This is feature variance - different concept.
			\item \textcolor{myred}{D:} This is correlation - different concept.
		\end{itemize}
		
		\textcolor{myred}{\textbf{Exam Trap:}} Feature importance helps \textcolor{myblue}{interpret} model predictions! This is the key purpose.
	\end{frame}
	
	\begin{frame}{Q98: Interpretability - SHAP Values}
		\fontsize{6.5}{7.5}\selectfont
		\textbf{What are SHAP (SHapley Additive exPlanations) values?}
		
		\begin{enumerate}
			\item A method to compute feature importance using game theory.
			\item A method to increase model accuracy.
			\item A method to regularize models.
			\item A method to select the best hyperparameters.
		\end{enumerate}
		
		\vspace{0.05cm}
		\textcolor{myblue}{\textbf{Answer: A}}
		\vspace{0.05cm}
		\textcolor{myblue}{\textbf{Explanation:}} SHAP values use Shapley values from game theory. Provide unified framework for interpreting model predictions by attributing prediction to each feature.
		
		\vspace{0.05cm}
		\textcolor{myred}{\textbf{Why Others Are Wrong:}}
		\begin{itemize}
			\item \textcolor{myred}{B:} SHAP doesn't increase accuracy - it's for explanation.
			\item \textcolor{myred}{C:} SHAP is not regularization.
			\item \textcolor{myred}{D:} SHAP is not hyperparameter tuning.
		\end{itemize}
		
		\textcolor{myred}{\textbf{Exam Trap:}} SHAP uses \textcolor{myblue}{game theory} to explain predictions! This is the key insight.
	\end{frame}
	
	\begin{frame}{Q99: Interpretability - LIME}
		\fontsize{6.5}{7.5}\selectfont
		\textbf{What is LIME (Local Interpretable Model-agnostic Explanations)?}
		
		\begin{enumerate}
			\item A method to explain individual predictions by approximating the model locally with an interpretable model.
			\item A method to increase model accuracy.
			\item A method to regularize models.
			\item A method to select the best features.
		\end{enumerate}
		
		\vspace{0.05cm}
		\textcolor{myblue}{\textbf{Answer: A}}
		\vspace{0.05cm}
		\textcolor{myblue}{\textbf{Explanation:}} LIME explains individual predictions by approximating the complex model locally (near a specific prediction) with an interpretable model (e.g., linear regression).
		
		\vspace{0.05cm}
		\textcolor{myred}{\textbf{Why Others Are Wrong:}}
		\begin{itemize}
			\item \textcolor{myred}{B:} LIME doesn't increase accuracy - it's for explanation.
			\item \textcolor{myred}{C:} LIME is not regularization.
			\item \textcolor{myred}{D:} LIME is not feature selection (though it identifies important features locally).
		\end{itemize}
		
		\textcolor{myred}{\textbf{Exam Trap:}} LIME provides \textcolor{myblue}{local explanations} for individual predictions! This is the key distinction from SHAP.
	\end{frame}
	
	\begin{frame}{Q100: Model Estimation - Final Summary}
		\fontsize{6.5}{7.5}\selectfont
		\textbf{Which of the following provides the most comprehensive summary of supervised learning model estimation?}
		
		\begin{enumerate}
			\item Only OLS is used for model estimation.
			\item Model estimation involves choosing between OLS, NLS, and MLE, optimizing using gradient descent, and evaluating with cross-validation.
			\item Only MLE is used for model estimation.
			\item Model estimation is only about regularization.
		\end{enumerate}
		
		\vspace{0.05cm}
		\textcolor{myblue}{\textbf{Answer: B}}
		\vspace{0.05cm}
		\textcolor{myblue}{\textbf{Explanation:}} Comprehensive estimation includes: choosing estimation method (OLS, NLS, MLE), optimizing (gradient descent), and evaluating (cross-validation).
		
		\vspace{0.05cm}
		\textcolor{myred}{\textbf{Why Others Are Wrong:}}
		\begin{itemize}
			\item \textcolor{myred}{A:} Only OLS is incomplete - NLS and MLE also used.
			\item \textcolor{myred}{C:} Only MLE is incomplete - OLS also used.
			\item \textcolor{myred}{D:} Regularization is part of model fitting, not the entire estimation process.
		\end{itemize}
		
		\textcolor{myred}{\textbf{Exam Trap:}} Model estimation is \textcolor{myblue}{comprehensive} - includes method choice, optimization, and evaluation! This is the big picture.
	\end{frame}
	
	% ============================================================
	% SUMMARY SLIDE
	% ============================================================
	
	\begin{frame}{Summary: Key Concepts}
		\fontsize{6.5}{7.5}\selectfont
		\textbf{Core Estimation Methods:}
		\begin{itemize}
			\item \textcolor{myblue}{OLS:} Minimizes SSE; closed-form for linear regression; BLUE under Gauss-Markov
			\item \textcolor{myblue}{MLE:} Maximizes likelihood; equivalent to OLS under normal errors
			\item \textcolor{myblue}{NLS:} For nonlinear-in-parameters models; requires numerical optimization
		\end{itemize}
		
		\vspace{0.1cm}
		\textbf{Optimization:}
		\begin{itemize}
			\item \textcolor{myblue}{Gradient Descent:} Iterative; uses negative gradient; learning rate controls step size
			\item \textcolor{myblue}{Backpropagation:} Chain rule for neural networks; requires storing activations
		\end{itemize}
		
		\vspace{0.1cm}
		\textbf{Model Selection:}
		\begin{itemize}
			\item \textcolor{myblue}{Bias-Variance Tradeoff:} Simplicity vs complexity
			\item \textcolor{myblue}{Regularization:} L1 (Lasso - feature selection), L2 (Ridge - shrinkage)
			\item \textcolor{myblue}{Cross-Validation:} k-fold (k=10 common), evaluates generalization
		\end{itemize}
		
		\textcolor{myred}{\textbf{Exam Success:}} Understand the tradeoffs and know the definitions!
	\end{frame}
	
\end{document}